\documentclass[12pt]{article}

% === Кодировки и язык ===
\usepackage{cmap}
\usepackage[T2A]{fontenc}
\usepackage{mathtext}
\usepackage[utf8]{inputenc}
\usepackage[english, russian]{babel}

% === Цвета ===
\usepackage[dvipsnames]{xcolor}

% === Математика ===
\usepackage{amsmath, amssymb, wasysym, amsthm}
\usepackage{stmaryrd}
\usepackage{proof}

% === Графика и float'ы ===
\usepackage{wrapfig, float}
\usepackage{graphicx}

% === TikZ / PGFPlots ===
\usepackage{pgfplots}
\pgfplotsset{compat=newest}
\usepackage{pgfplotstable}
\usepackage{tikz, tikz-3dplot}
\usetikzlibrary{calc, patterns, angles, quotes, intersections, automata, positioning, arrows.meta}

% === Таблицы ===
\usepackage{tabularx}
\usepackage{xltabular}
\renewcommand\tabularxcolumn[1]{m{#1}}

% === Разметка ===
\usepackage{geometry}
\geometry{margin=0.5in}

% === Списки и прочее ===
\usepackage[shortlabels]{enumitem}
\usepackage{verbatim}

% ==========================================
% === ТЕОРЕМНЫЕ ОКРУЖЕНИЯ (amsthm) ===
% ==========================================
\theoremstyle{plain}
\newtheorem{theorem}{Теорема}[section]
\newtheorem{corollary}{Следствие}[section]
\newtheorem{lemma}{Лемма}[section]
\newtheorem{statement}{Утверждение}[section]
\newtheorem{axiom}{Аксиома}[section]
\newtheorem{law}{Закон}[section]

\theoremstyle{definition}
\newtheorem{definition}{Определение}[section]
\newtheorem{example}{Пример}[section]
\newtheorem{problem}{Задача}
\newtheorem{method}{Метод}[section]
\newtheorem{event}{Событие}[section]

\theoremstyle{remark}
\newtheorem{note}{Заметка}[section]
\newtheorem{property}{Свойство}[section]
\newtheorem{properties}{Свойства}[section]
\newtheorem{principle}{Правило}[section]
\newtheorem{person}{Личность}[section]
\newtheorem*{solution}{Решение}

% ==========================================
% === КРАСИВЫЕ РАМКИ (tcolorbox) ===
% ==========================================
\usepackage{tcolorbox}
\tcbuselibrary{skins, breakable, theorems}

% Базовые стили для разных типов блоков
\tcbset{
    % Стиль для теорем, лемм (Синий - строгий)
    thmstyle/.style={
        enhanced, breakable,
        frame hidden, boxrule=0pt,
        colback=blue!5!white, % Фон: 5% синего, 95% белого
        borderline west={2.5pt}{0pt}{blue!70!black}, % Вертикальная линия слева
        before skip=10pt, after skip=10pt,
        left=4pt, right=4pt, top=4pt, bottom=4pt
    },
    % Стиль для доказательств (фиолетовый - процесс рассуждения)
    proofstyle/.style={
        enhanced, breakable,
        frame hidden, boxrule=0pt,
        colback=Plum!5!white,
        borderline west={2.5pt}{0pt}{Plum!80!black},
        before skip=10pt, after skip=10pt,
        left=4pt, right=4pt, top=4pt, bottom=4pt,
        sharp corners
    },
    % Стиль для определений и аксиом (Зеленый - фундаментальный)
    defstyle/.style={
        enhanced, breakable,
        frame hidden, boxrule=0pt,
        colback=ForestGreen!5!white,
        borderline west={2.5pt}{0pt}{ForestGreen!80!black},
        before skip=10pt, after skip=10pt,
        left=4pt, right=4pt, top=4pt, bottom=4pt
    },
    % Стиль для примеров и задач (Оранжевый - практический)
    exstyle/.style={
        enhanced, breakable,
        frame hidden, boxrule=0pt,
        colback=orange!5!white,
        borderline west={2.5pt}{0pt}{orange!90!black},
        before skip=10pt, after skip=10pt,
        left=4pt, right=4pt, top=4pt, bottom=4pt
    },
    % Стиль для заметок и свойств (Серый - нейтральный)
    remstyle/.style={
        enhanced, breakable,
        frame hidden, boxrule=0pt,
        colback=gray!5!white,
        borderline west={2.5pt}{0pt}{gray!70!black},
        before skip=10pt, after skip=10pt,
        left=4pt, right=4pt, top=4pt, bottom=4pt
    }
}

% Применяем стили к существующим окружениям amsthm
% 1. Строгая математика
\tcolorboxenvironment{theorem}{thmstyle}
\tcolorboxenvironment{lemma}{thmstyle}
\tcolorboxenvironment{corollary}{thmstyle}
\tcolorboxenvironment{statement}{thmstyle}
\tcolorboxenvironment{law}{thmstyle}

% 2. Доказательство
\tcolorboxenvironment{proof}{proofstyle}
\renewcommand{\qed}{}

% 3. Определения
\tcolorboxenvironment{definition}{defstyle}
\tcolorboxenvironment{axiom}{defstyle}

% 4. Практика
\tcolorboxenvironment{example}{exstyle}
\tcolorboxenvironment{problem}{exstyle}
\tcolorboxenvironment{method}{exstyle}
\tcolorboxenvironment{event}{exstyle}

% 5. Заметки и прочее
\tcolorboxenvironment{note}{remstyle}
\tcolorboxenvironment{property}{remstyle}
\tcolorboxenvironment{properties}{remstyle}
\tcolorboxenvironment{principle}{remstyle}
\tcolorboxenvironment{person}{remstyle}
\tcolorboxenvironment{solution}{remstyle}

% === Гиперссылки (почти последним!) ===
\usepackage[
colorlinks=true,
linkcolor=blue,
urlcolor=blue,
citecolor=blue,
bookmarks=true,
bookmarksopen=false
]{hyperref}


\newcommand{\Cl}{\text{Cl}}

\begin{document}
	\tableofcontents
	\setcounter{tocdepth}{3}
	\newpage
	
	\section{Числа Фибоначчи. Задачи, приводящие к числам Фибоначчи. Некоторые тождества с числами Фибоначчи. НОД чисел Фибоначчи. Зацикливание остатков и поведение чисел Фибоначчи, кратных n.}
	
	\subsection{Определение и задачи}
	
	\begin{definition}
		Числа Фибоначчи определяются рекуррентно: $F_1 = F_2 = 1$, $F_{n+1} = F_n + F_{n-1}$ для $n \geqslant 2$.
	\end{definition}
	
	Первые несколько чисел: $1, 1, 2, 3, 5, 8, 13, 21, 34, 55, \dots$
	
	Удобно также доопределить $F_0 = 0$ (тогда рекуррентность $F_{n+1} = F_n + F_{n-1}$ выполняется и при $n = 1$).
	
	\textbf{Задачи, приводящие к числам Фибоначчи:}
	\begin{enumerate}
		\item \textbf{Задача о кроликах (Фибоначчи, 1202):} Сколько пар кроликов будет через $n$ месяцев, если каждая пара начинает размножаться со второго месяца жизни и производит одну пару в месяц?
		
		\textit{Решение:} Пусть $a_n$~--- число пар на $n$-м месяце. Пары делятся на «взрослых» (которые уже были в прошлом месяце, их $a_{n-1}$) и «новорождённых» (их столько, сколько взрослых пар в прошлом месяце, то есть $a_{n-2}$). Итого $a_n = a_{n-1} + a_{n-2}$ при $a_1 = a_2 = 1$, то есть $a_n = F_n$.
		
		\item \textbf{Задача о лестнице:} Сколькими способами можно подняться на $n$-ю ступеньку, если можно шагать на 1 или 2 ступеньки?
		
		\textit{Решение:} Пусть $a_n$~--- число способов. Последний шаг был либо с $(n-1)$-й ступеньки (шаг на 1), либо с $(n-2)$-й (шаг на 2). Значит $a_n = a_{n-1} + a_{n-2}$ при $a_1 = 1$, $a_2 = 2$, то есть $a_n = F_{n+1}$.
		
		\item \textbf{Задача о замощении:} Число способов замостить полоску $2 \times n$ плитками $1 \times 2$ и $2 \times 1$ равно $F_{n+1}$.
		
		\textit{Решение:} Пусть $a_n$~--- число способов замощения $2 \times n$. Смотрим на левый край: если там стоит вертикальная плитка, остаток~--- $2 \times (n-1)$; если две горизонтальных (друг над другом), остаток~--- $2 \times (n-2)$. Значит $a_n = a_{n-1} + a_{n-2}$ при $a_1 = 1$, $a_2 = 2$, откуда $a_n = F_{n+1}$.
	\end{enumerate}
	
	\subsection{Свойства чисел Фибоначчи}
	
	\begin{enumerate}
		\item $F_1 + F_2 + \dots + F_n = F_{n + 2} - 1$
		
		\begin{proof}
			Индукция по $n$. База: $n = 1$: $F_1 = 1 = F_3 - 1 = 2 - 1$. \\
			Переход: пусть верно для $n$. Тогда
			\[F_1 + \dots + F_n + F_{n+1} = (F_{n+2} - 1) + F_{n+1} = F_{n+3} - 1.\]
		\end{proof}
		
		\item $F_1 + F_3 + \dots + F_{2n - 1} = F_{2n}$
		
		\begin{proof}
			Индукция по $n$. База: $n = 1$: $F_1 = 1 = F_2$. \\
			Переход: $F_1 + F_3 + \dots + F_{2n-1} + F_{2n+1} = F_{2n} + F_{2n+1} = F_{2n+2}$.
		\end{proof}
		
		\item $F_2 + F_4 + \dots + F_{2n} = F_{2n + 1} - 1$
		
		\begin{proof}
			Из свойств 1 и 2: 
			\begin{align*}
				F_2 + F_4 + \dots + F_{2n} &= (F_1 + F_2 + \dots + F_{2n}) - (F_1 + F_3 + \dots + F_{2n-1}) \\
				&= (F_{2n+2} - 1) - F_{2n} = F_{2n+1} - 1.
			\end{align*}
		\end{proof}
		
		\item $F_1^2 + F_2^2 + \dots + F_n^2 = F_n F_{n + 1}$
		
		\begin{proof}
			Индукция по $n$. База: $n = 1$: $F_1^2 = 1 = F_1 F_2 = 1 \cdot 1$. \\
			Переход: 
			\[F_1^2 + \dots + F_n^2 + F_{n+1}^2 = F_n F_{n+1} + F_{n+1}^2 = F_{n+1}(F_n + F_{n+1}) = F_{n+1} F_{n+2}.\]
		\end{proof}
		
		\item (Тождество Кассини) $F_{n - 1} F_{n + 1} - F_n^2 = (-1)^n$
		
		\begin{proof}
			Индукция по $n$. База: $n = 2$: $F_1 F_3 - F_2^2 = 1 \cdot 2 - 1 = 1 = (-1)^2$. \\
			Переход: пусть $F_{n-1} F_{n+1} - F_n^2 = (-1)^n$. Тогда
			\begin{align*}
				F_n F_{n+2} - F_{n+1}^2 &= F_n(F_n + F_{n+1}) - F_{n+1}^2 \\
				&= F_n^2 + F_n F_{n+1} - F_{n+1}^2 \\
				&= F_n^2 - F_{n+1}(F_{n+1} - F_n) \\
				&= F_n^2 - F_{n+1} F_{n-1} = -(-1)^n = (-1)^{n+1}.
			\end{align*}
		\end{proof}
		
		\begin{note}
			Матричная интерпретация: $\begin{pmatrix} F_{n+1} & F_n \\ F_n & F_{n-1} \end{pmatrix} = \begin{pmatrix} 1 & 1 \\ 1 & 0 \end{pmatrix}^n$. Тождество Кассини~--- это равенство определителей: $\det\begin{pmatrix} 1 & 1 \\ 1 & 0 \end{pmatrix}^n = \left(\det\begin{pmatrix} 1 & 1 \\ 1 & 0 \end{pmatrix}\right)^n = (-1)^n$.
		\end{note}
		
		\item $F_{n + k} = F_n F_{k - 1} + F_{n + 1} F_k$
		
		\begin{proof}
			Индукция по $k$ при фиксированном $n$. \\
			База: $k = 1$: $F_{n+1} = F_n \cdot F_0 + F_{n+1} \cdot F_1 = F_n \cdot 0 + F_{n+1} \cdot 1 = F_{n+1}$. \\
			$k = 2$: $F_{n+2} = F_n \cdot F_1 + F_{n+1} \cdot F_2 = F_n + F_{n+1} = F_{n+2}$. \\
			Переход: пусть верно для $k$ и $k-1$. Тогда
			\begin{align*}
				F_{n+k+1} &= F_{n+k} + F_{n+k-1} \\
				&= (F_n F_{k-1} + F_{n+1} F_k) + (F_n F_{k-2} + F_{n+1} F_{k-1}) \\
				&= F_n(F_{k-1} + F_{k-2}) + F_{n+1}(F_k + F_{k-1}) \\
				&= F_n F_k + F_{n+1} F_{k+1}.
			\end{align*}
		\end{proof}
		
		\begin{note}
			Частный случай $n = k$: $F_{2n} = F_n F_{n-1} + F_{n+1} F_n = F_n(F_{n-1} + F_{n+1})$, что даёт удобную формулу для быстрого вычисления чисел Фибоначчи.
		\end{note}
	\end{enumerate}
	
	\subsection{НОД чисел Фибоначчи}
	
	\begin{theorem}
		$(F_m, F_n) = F_{(m, n)}$
	\end{theorem}
	
	\begin{proof}
		\textbf{Шаг 1:} Докажем, что $(F_m, F_n) = (F_{m-n}, F_n)$ при $m > n$.
		
		Из свойства 6 (с заменой $k \to n$, $n \to m - n$): $F_m = F_{m-n} F_{n-1} + F_{m-n+1} F_n$.
		
		Значит, $F_{m-n} F_{n-1} = F_m - F_{m-n+1} F_n$, откуда:
		\[(F_m, F_n) = (F_{m-n} F_{n-1} + F_{m-n+1} F_n, F_n) = (F_{m-n} F_{n-1}, F_n).\]
		
		(Последнее равенство: $\gcd(a + bc, c) = \gcd(a, c)$.)
		
		Покажем, что $(F_{n-1}, F_n) = 1$. По тождеству Кассини: $F_{n-1} F_{n+1} - F_n^2 = (-1)^n$.
		
		Если $d = (F_{n-1}, F_n)$, то $d \mid F_{n-1}$ и $d \mid F_n$, значит $d \mid F_{n-1} F_{n+1}$ и $d \mid F_n^2$, откуда $d \mid (F_{n-1} F_{n+1} - F_n^2) = (-1)^n$, то есть $d = 1$.
		
		Следовательно, $(F_{m-n} F_{n-1}, F_n) = (F_{m-n}, F_n)$~--- поскольку $F_{n-1}$ взаимно просто с $F_n$, общий делитель $F_{m-n}$ и $F_n$ не «улучшается» множителем $F_{n-1}$.
		
		\textbf{Шаг 2:} По аналогии с алгоритмом Евклида, повторяя вычитание:
		\[(F_m, F_n) = (F_{m-n}, F_n) = (F_{m-2n}, F_n) = \dots = (F_{m \bmod n}, F_n).\]
		
		Затем меняем роли слагаемых и продолжаем, как в алгоритме Евклида для $m$ и $n$:
		\[(F_m, F_n) \to (F_{m \bmod n}, F_n) \to \dots \to (F_{(m,n)}, F_{(m,n)}) = F_{(m,n)}.\]
		
		Последнее равенство: $(F_k, F_k) = F_k$.
	\end{proof}
	
	\begin{corollary}
		$F_m \mid F_n \Leftrightarrow m \mid n$.
	\end{corollary}
	
	\begin{proof}
		$F_m \mid F_n \Leftrightarrow (F_m, F_n) = F_m \Leftrightarrow F_{(m,n)} = F_m \Leftrightarrow (m, n) = m \Leftrightarrow m \mid n$.
	\end{proof}
	
	\subsection{Зацикливание остатков}
	
	\begin{theorem}
		Последовательность $F_n \bmod m$ периодична для любого $m \geqslant 2$.
	\end{theorem}
	
	\begin{proof}
		Рассмотрим пары $(F_n \bmod m, F_{n+1} \bmod m)$. Всего различных пар не более $m^2$ (каждая компонента принимает значения в $\{0, 1, \dots, m-1\}$).
		
		По принципу Дирихле, среди первых $m^2 + 1$ пар найдутся две одинаковые: $(F_i \bmod m, F_{i+1} \bmod m) = (F_j \bmod m, F_{j+1} \bmod m)$ при $i < j$.
		
		Поскольку каждый следующий член однозначно определяется двумя предыдущими по рекуррентности $F_{n+2} \equiv F_{n+1} + F_n \pmod{m}$, имеем $F_{i+k} \equiv F_{j+k} \pmod{m}$ для всех $k \geqslant 0$, то есть последовательность является периодической начиная с позиции $i$.
		
		Кроме того, рекуррентность работает и назад: $F_{n-1} = F_{n+1} - F_n \equiv F_{n+1} - F_n \pmod{m}$. Значит из совпадения пары $(F_i, F_{i+1})$ с $(F_j, F_{j+1})$ следует совпадение $(F_{i-1}, F_i)$ с $(F_{j-1}, F_j)$, и так далее. Идя назад до позиции $n = 1$, получаем, что последовательность периодична с периодом $j - i$ уже с самого начала.
	\end{proof}
	
	\begin{definition}
		Период Пизано $\pi(m)$~--- наименьший период последовательности $F_n \bmod m$.
	\end{definition}
	
	\begin{example}
		$\pi(2) = 3$: $F_n \bmod 2 = 1, 1, 0, 1, 1, 0, \dots$
		
		$\pi(3) = 8$: $F_n \bmod 3 = 1, 1, 2, 0, 2, 2, 1, 0, 1, 1, \dots$
		
		$\pi(5) = 20$, $\pi(10) = 60$.
	\end{example}
	
	\begin{theorem}
		Числа Фибоначчи, кратные $n$, появляются периодически: $n \mid F_k \Leftrightarrow \alpha(n) \mid k$, где $\alpha(n)$~--- наименьший индекс такой, что $n \mid F_{\alpha(n)}$ (ранг вхождения). Этот период $\alpha(n)$ делит $\pi(n)$.
	\end{theorem}
	
	\begin{proof}
		Из формулы $F_{n+k} = F_n F_{k-1} + F_{n+1} F_k$ при $n = \alpha(n)$:
		$F_{\alpha(n) + k} = F_{\alpha(n)} F_{k-1} + F_{\alpha(n)+1} F_k \equiv 0 \cdot F_{k-1} + F_{\alpha(n)+1} F_k \pmod{n}$.
		
		При $k = \alpha(n)$: $F_{2\alpha(n)} \equiv F_{\alpha(n)+1} F_{\alpha(n)} \equiv 0 \pmod{n}$.
		
		Индукцией доказывается $F_{t \cdot \alpha(n)} \equiv 0 \pmod{n}$ для всех $t \geqslant 0$.
		
		Обратно, если $k$ не кратно $\alpha(n)$, то $F_k \not\equiv 0 \pmod{n}$ (это следует из того, что $(F_m, F_n) = F_{(m,n)}$, поэтому $n \mid F_k$ влечёт $\alpha(n) \mid k$).
		
		Наконец, $\alpha(n) \mid \pi(n)$, так как $F_{\pi(n)} \equiv F_0 = 0 \pmod{n}$ (последовательность периодична с периодом $\pi(n)$, а $F_0 = 0$).
	\end{proof}
	
	\newpage
	
	\section{Множества, операции над множествами. Отображения: сюръекции, инъекции, биекции. Образы и прообразы множеств. Односторонне и двусторонне обратные отображения. Отношения, отношения эквивалентности, отношения порядка.}
	
	\subsection{Множества и операции над ними}
	
	\begin{definition}
		Множество~--- совокупность различных объектов, рассматриваемых как единое целое.
	\end{definition}
	
	\textbf{Операции над множествами:}
	\begin{itemize}
		\item \textbf{Объединение:} $A \cup B = \{x : x \in A \text{ или } x \in B\}$
		\item \textbf{Пересечение:} $A \cap B = \{x : x \in A \text{ и } x \in B\}$
		\item \textbf{Разность:} $A \setminus B = \{x : x \in A \text{ и } x \notin B\}$
		\item \textbf{Симметрическая разность:} $A \triangle B = (A \setminus B) \cup (B \setminus A)$
		\item \textbf{Дополнение:} $\overline{A} = X \setminus A$ (относительно универсума $X$)
		\item \textbf{Декартово произведение:} $A \times B = \{(a, b) : a \in A, b \in B\}$
	\end{itemize}
	
	\textbf{Законы де Моргана:}
	\[\overline{A \cup B} = \overline{A} \cap \overline{B}, \quad \overline{A \cap B} = \overline{A} \cup \overline{B}\]
	
	\begin{proof}
		Докажем первый закон. $x \in \overline{A \cup B} \Leftrightarrow x \notin A \cup B \Leftrightarrow (x \notin A) \text{ и } (x \notin B) \Leftrightarrow x \in \overline{A} \text{ и } x \in \overline{B} \Leftrightarrow x \in \overline{A} \cap \overline{B}$. Второй аналогично.
	\end{proof}
	
	\subsection{Отображения}
	
	\begin{definition}
		Отображение $f: X \to Y$~--- правило, сопоставляющее каждому $x \in X$ единственный элемент $f(x) \in Y$. Множество $X$ называется \textbf{областью определения}, $Y$~--- \textbf{областью значений (кодомен)}.
	\end{definition}
	
	\begin{definition}
		Отображение $f: X \to Y$ называется:
		\begin{itemize}
			\item \textbf{Инъекцией} (вложением), если $f(x_1) = f(x_2) \Rightarrow x_1 = x_2$, то есть различные элементы отображаются в различные.
			\item \textbf{Сюръекцией} (отображением «на»), если $\forall y \in Y \ \exists x \in X: f(x) = y$, то есть каждый элемент кодомена имеет хотя бы один прообраз.
			\item \textbf{Биекцией} (взаимно однозначным соответствием), если $f$ инъективно и сюръективно, то есть каждый элемент $y \in Y$ имеет ровно один прообраз.
		\end{itemize}
	\end{definition}
	
	\begin{note}
		Для конечных множеств $|X| = |Y|$: инъекция $\Leftrightarrow$ сюръекция $\Leftrightarrow$ биекция.
	\end{note}
	
	\begin{definition}
		\textbf{Изоморфизм}~--- биекция, сохраняющая структуру. В зависимости от контекста:
		\begin{itemize}
			\item \textbf{Изоморфизм частично упорядоченных множеств:} биекция $f: (X, \leq_X) \to (Y, \leq_Y)$ такая, что $x_1 \leq_X x_2 \Leftrightarrow f(x_1) \leq_Y f(x_2)$.
			\item \textbf{Изоморфизм групп:} биективный гомоморфизм $f: G \to H$, то есть $f(ab) = f(a)f(b)$.
			\item \textbf{Изоморфизм графов:} биекция $f: V(G) \to V(H)$, сохраняющая смежность: $\{u, v\} \in E(G) \Leftrightarrow \{f(u), f(v)\} \in E(H)$.
			\item \textbf{Изометрия метрических пространств:} биекция $f: (X, d_X) \to (Y, d_Y)$ такая, что $d_Y(f(x), f(y)) = d_X(x, y)$.
		\end{itemize}
	\end{definition}
	
	\begin{note}
		Два объекта называются \textbf{изоморфными} ($X \cong Y$), если между ними существует изоморфизм. Изоморфные объекты неразличимы с точки зрения соответствующей структуры.
	\end{note}
	
	\begin{theorem}
		Изоморфизм~--- отношение эквивалентности:
		\begin{enumerate}
			\item (Рефлексивность) $X \cong X$ (тождественное отображение).
			\item (Симметричность) $X \cong Y \Rightarrow Y \cong X$ (обратное отображение).
			\item (Транзитивность) $X \cong Y, Y \cong Z \Rightarrow X \cong Z$ (композиция).
		\end{enumerate}
	\end{theorem}
	
	\begin{proof}
		1. Тождественное отображение $\text{id}: X \to X$~--- биекция, сохраняющая структуру.
		
		2. Если $f: X \to Y$~--- изоморфизм, то $f^{-1}: Y \to X$~--- тоже биекция.
		
		Сохранение структуры: например, для частичных порядков: $y_1 \leq_Y y_2 \Leftrightarrow f(f^{-1}(y_1)) \leq_Y f(f^{-1}(y_2)) \Leftrightarrow f^{-1}(y_1) \leq_X f^{-1}(y_2)$.
		
		Для групп: $f^{-1}(ab) = f^{-1}(f(f^{-1}(a)) \cdot f(f^{-1}(b))) = f^{-1}(f(f^{-1}(a) \cdot f^{-1}(b))) = f^{-1}(a) \cdot f^{-1}(b)$.
		
		3. Если $f: X \to Y$ и $g: Y \to Z$~--- изоморфизмы, то $g \circ f: X \to Z$~--- биекция.
		
		Сохранение структуры: например, для групп: $(g \circ f)(ab) = g(f(ab)) = g(f(a)f(b)) = g(f(a))g(f(b)) = (g \circ f)(a) \cdot (g \circ f)(b)$.
	\end{proof}
	
	\subsection{Образы и прообразы}
	
	\begin{definition}
		Пусть $f: X \to Y$, $A \subseteq X$, $B \subseteq Y$. Тогда:
		\begin{itemize}
			\item \textbf{Образ:} $f(A) = \{f(x) : x \in A\} \subseteq Y$
			\item \textbf{Прообраз:} $f^{-1}(B) = \{x \in X : f(x) \in B\} \subseteq X$
		\end{itemize}
	\end{definition}
	
	\begin{note}
		Запись $f^{-1}(B)$ обозначает прообраз множества и не предполагает существования обратного отображения $f^{-1}$.
	\end{note}
	
	\begin{theorem}[Свойства образов и прообразов]
		\begin{enumerate}
			\item $f^{-1}(B_1 \cup B_2) = f^{-1}(B_1) \cup f^{-1}(B_2)$
			\item $f^{-1}(B_1 \cap B_2) = f^{-1}(B_1) \cap f^{-1}(B_2)$
			\item $f^{-1}(Y \setminus B) = X \setminus f^{-1}(B)$
			\item $f(A_1 \cup A_2) = f(A_1) \cup f(A_2)$
			\item $f(A_1 \cap A_2) \subseteq f(A_1) \cap f(A_2)$ (равенство при инъективности)
			\item $A \subseteq f^{-1}(f(A))$ (равенство при инъективности)
			\item $f(f^{-1}(B)) \subseteq B$ (равенство при сюръективности)
		\end{enumerate}
	\end{theorem}
	
	\begin{proof}
		1. $x \in f^{-1}(B_1 \cup B_2) \Leftrightarrow f(x) \in B_1 \cup B_2 \Leftrightarrow f(x) \in B_1 \text{ или } f(x) \in B_2 \Leftrightarrow x \in f^{-1}(B_1) \cup f^{-1}(B_2)$.
		
		2. $x \in f^{-1}(B_1 \cap B_2) \Leftrightarrow f(x) \in B_1 \text{ и } f(x) \in B_2 \Leftrightarrow x \in f^{-1}(B_1) \cap f^{-1}(B_2)$.
		
		3. $x \in f^{-1}(Y \setminus B) \Leftrightarrow f(x) \notin B \Leftrightarrow x \notin f^{-1}(B) \Leftrightarrow x \in X \setminus f^{-1}(B)$.
		
		4. $y \in f(A_1 \cup A_2) \Leftrightarrow \exists x \in A_1 \cup A_2: f(x) = y \Leftrightarrow y \in f(A_1) \cup f(A_2)$.
		
		5. Если $y \in f(A_1 \cap A_2)$, то $\exists x \in A_1 \cap A_2: f(x) = y$, значит $y \in f(A_1)$ и $y \in f(A_2)$, то есть $y \in f(A_1) \cap f(A_2)$. Обратное включение неверно в общем случае: если $f$ не инъективна, возможно $y = f(x_1) = f(x_2)$, $x_1 \in A_1 \setminus A_2$, $x_2 \in A_2 \setminus A_1$~--- тогда $y \in f(A_1) \cap f(A_2)$, но $y \notin f(A_1 \cap A_2)$. При инъективности это невозможно.
		
		6. Если $x \in A$, то $f(x) \in f(A)$, значит $x \in f^{-1}(f(A))$. При инъективности: если $x \in f^{-1}(f(A))$, то $f(x) \in f(A)$, то есть $\exists x' \in A: f(x') = f(x)$; по инъективности $x = x' \in A$.
		
		7. Если $y \in f(f^{-1}(B))$, то $\exists x \in f^{-1}(B): f(x) = y$, то есть $f(x) \in B$, значит $y \in B$. При сюръективности: для любого $y \in B$ существует $x: f(x) = y$, тогда $x \in f^{-1}(B)$ и $y = f(x) \in f(f^{-1}(B))$.
	\end{proof}
	
	\subsection{Обратные отображения}
	
	\begin{definition}
		Пусть $f: X \to Y$.
		\begin{itemize}
			\item $g: Y \to X$~--- \textbf{левое обратное}, если $g \circ f = \text{id}_X$ (то есть $g(f(x)) = x$ для всех $x \in X$).
			\item $h: Y \to X$~--- \textbf{правое обратное}, если $f \circ h = \text{id}_Y$ (то есть $f(h(y)) = y$ для всех $y \in Y$).
			\item $f^{-1}: Y \to X$~--- \textbf{двустороннее обратное}, если $f^{-1} \circ f = \text{id}_X$ и $f \circ f^{-1} = \text{id}_Y$.
		\end{itemize}
	\end{definition}
	
	\begin{theorem}
		\begin{enumerate}
			\item $f$ имеет левое обратное $\Leftrightarrow$ $f$ инъективно
			\item $f$ имеет правое обратное $\Leftrightarrow$ $f$ сюръективно
			\item $f$ имеет двустороннее обратное $\Leftrightarrow$ $f$ биективно
		\end{enumerate}
	\end{theorem}
	
	\begin{proof}
		1. $(\Rightarrow)$ Пусть $g \circ f = \text{id}_X$. Если $f(x_1) = f(x_2)$, то $x_1 = g(f(x_1)) = g(f(x_2)) = x_2$.
		
		$(\Leftarrow)$ Пусть $f$ инъективно. Зафиксируем $x_0 \in X$. Определим $g: Y \to X$ по правилу: $g(y) = x$, если $y = f(x)$ для некоторого (единственного, по инъективности) $x \in X$; $g(y) = x_0$, если $y \notin f(X)$. Тогда $g(f(x)) = x$ для всех $x$.
		
		2. $(\Rightarrow)$ Пусть $f \circ h = \text{id}_Y$. Для любого $y \in Y$: $f(h(y)) = y$, значит $y \in f(X)$.
		
		$(\Leftarrow)$ Пусть $f$ сюръективно. Для каждого $y \in Y$ выберем $h(y) \in X$ такой, что $f(h(y)) = y$ (такой элемент существует по сюръективности; выбор требует аксиомы выбора в общем случае). Тогда $f(h(y)) = y$.
		
		3. $(\Rightarrow)$ Если $f$ имеет двустороннее обратное $f^{-1}$, то $f^{-1}$~--- одновременно левое и правое, значит $f$ инъективно (из 1) и сюръективно (из 2).
		
		$(\Leftarrow)$ Если $f$ биективно, то оно имеет левое обратное $g$ (из 1) и правое $h$ (из 2). Покажем, что $g = h$:
		\[g = g \circ \text{id}_Y = g \circ (f \circ h) = (g \circ f) \circ h = \text{id}_X \circ h = h.\]
		
		Значит $g = h$~--- двустороннее обратное к $f$.
	\end{proof}
	
	\begin{note}
		Двустороннее обратное единственно: если $g$ и $g'$~--- оба двусторонних обратных, то $g = g \circ \text{id} = g \circ (f \circ g') = (g \circ f) \circ g' = \text{id} \circ g' = g'$.
	\end{note}
	
	\subsection{Отношения}
	
	\begin{definition}
		Бинарное отношение на $X$~--- подмножество $R \subseteq X \times X$. Пишем $xRy$ вместо $(x, y) \in R$.
	\end{definition}
	
	\begin{definition}
		Отношение $\sim$ на $X$~--- \textbf{отношение эквивалентности}, если:
		\begin{itemize}
			\item (Рефлексивность) $x \sim x$ для всех $x \in X$
			\item (Симметричность) $x \sim y \Rightarrow y \sim x$
			\item (Транзитивность) $x \sim y, y \sim z \Rightarrow x \sim z$
		\end{itemize}
	\end{definition}
	
	\begin{definition}
		Класс эквивалентности элемента $x$: $[x] = \{y \in X : y \sim x\}$.
	\end{definition}
	
	\begin{theorem}
		Классы эквивалентности образуют разбиение множества $X$, то есть:
		\begin{enumerate}
			\item Каждый $x \in X$ принадлежит хотя бы одному классу.
			\item Любые два класса либо совпадают, либо не пересекаются.
		\end{enumerate}
	\end{theorem}
	
	\begin{proof}
		1. $x \in [x]$ по рефлексивности.
		
		2. Пусть $[x] \cap [y] \neq \varnothing$, возьмём $z \in [x] \cap [y]$. Тогда $z \sim x$ и $z \sim y$.
		
		По симметричности: $x \sim z$. По транзитивности: $x \sim y$.
		
		Покажем $[x] = [y]$: если $w \in [x]$, то $w \sim x \sim y$, значит $w \in [y]$, то есть $[x] \subseteq [y]$. Аналогично $[y] \subseteq [x]$.
	\end{proof}
	
	\begin{definition}
		\textbf{Фактор-множество} $X / {\sim}$~--- множество всех классов эквивалентности:
		\[X / {\sim} \; = \{[x] : x \in X\}.\]
	\end{definition}
	
	\begin{example}
		\begin{enumerate}
			\item $\mathbb{Z} / {\sim}$, где $a \sim b \Leftrightarrow a \equiv b \pmod{n}$~--- это $\mathbb{Z}_n = \{[0], [1], \dots, [n-1]\}$.
			\item $\mathbb{R} / {\sim}$, где $x \sim y \Leftrightarrow x - y \in \mathbb{Z}$~--- это «круг» $S^1$ (каждый класс~--- это набор чисел, отличающихся на целое).
			\item Пусть $f: X \to Y$~--- произвольное отображение. Тогда $x \sim y \Leftrightarrow f(x) = f(y)$~--- отношение эквивалентности, и $X / {\sim}$ биективно $\operatorname{Im} f$.
		\end{enumerate}
	\end{example}
	
	\begin{theorem}[Универсальное свойство фактор-множества]
		Пусть $\sim$~--- отношение эквивалентности на $X$, $\pi: X \to X/{\sim}$~--- естественная проекция $\pi(x) = [x]$.
		Тогда для любого отображения $f: X \to Y$, постоянного на классах эквивалентности (то есть $x \sim y \Rightarrow f(x) = f(y)$), существует единственное отображение $\tilde{f}: X/{\sim} \to Y$ такое, что $f = \tilde{f} \circ \pi$.
	\end{theorem}
	
	\begin{proof}
		\textbf{Существование:} определим $\tilde{f}([x]) = f(x)$.
		
		Корректность: если $[x] = [y]$, то $x \sim y$, значит $f(x) = f(y)$, то есть $\tilde{f}([x]) = \tilde{f}([y])$.
		
		Тогда $(\tilde{f} \circ \pi)(x) = \tilde{f}([x]) = f(x)$, то есть $\tilde{f} \circ \pi = f$.
		
		\textbf{Единственность:} если $\tilde{f} \circ \pi = f$, то $\tilde{f}([x]) = \tilde{f}(\pi(x)) = f(x)$~--- значение однозначно определено на каждом классе.
	\end{proof}
	
	\begin{definition}
		Отношение $\leqslant$ на $X$~--- \textbf{частичный порядок}, если:
		\begin{itemize}
			\item (Рефлексивность) $x \leqslant x$
			\item (Антисимметричность) $x \leqslant y, y \leqslant x \Rightarrow x = y$
			\item (Транзитивность) $x \leqslant y, y \leqslant z \Rightarrow x \leqslant z$
		\end{itemize}
		Пара $(X, \leqslant)$ называется \textbf{частично упорядоченным множеством (ч.у.м.)}.
	\end{definition}
	
	\begin{definition}
		Частичный порядок~--- \textbf{линейный (полный, тотальный)}, если $\forall x, y$: $x \leqslant y$ или $y \leqslant x$.
	\end{definition}
	
	\begin{example}
		\begin{itemize}
			\item $(\mathbb{N}, \mid)$~--- частичный, но не линейный порядок (2 и 3 несравнимы).
			\item $(\mathbb{R}, \leqslant)$~--- линейный порядок.
			\item $(2^X, \subseteq)$~--- частичный порядок для $|X| \geqslant 2$.
		\end{itemize}
	\end{example}
	
	\subsection{Теоремы Дилуорса и Мирского}
	
	\begin{definition}
		Пусть $(X, \leq)$~--- частично упорядоченное множество.
		\begin{itemize}
			\item \textbf{Цепь}~--- подмножество, в котором любые два элемента сравнимы.
			\item \textbf{Антицепь}~--- подмножество, в котором никакие два элемента не сравнимы.
		\end{itemize}
	\end{definition}
	
	\begin{theorem}[Дилуорс]
		Пусть $(X, \leq)$~--- конечное частично упорядоченное множество, и максимальный размер антицепи равен $k$. Тогда $X$ можно разбить на $k$ цепей, и не менее.
	\end{theorem}
	
	\begin{proof}
		Не менее $k$ цепей необходимо: каждая цепь содержит не более одного элемента любой антицепи, поэтому для покрытия антицепи размера $k$ нужно не менее $k$ цепей.
		
		Докажем, что $k$ цепей достаточно. Индукция по $|X|$. База $|X| = 1$ очевидна.
		
		Переход. Пусть максимальная антицепь имеет размер $k$.
		
		\textbf{Случай 1:} существует максимальная антицепь $A$ размера $k$, которая не совпадает ни со множеством всех максимальных, ни со множеством всех минимальных элементов.
		
		Определим:
		\[X^+ = \{x \in X : x \geq a \text{ для некоторого } a \in A\}, \quad X^- = \{x \in X : x \leq a \text{ для некоторого } a \in A\}.\]
		
		Ясно, что $A \subseteq X^+ \cap X^-$.
		
		\textit{Почему $X^+ \cup X^- = X$:} если $x \notin X^+ \cup X^-$, то $x$ несравним с каждым элементом $A$, значит $A \cup \{x\}$~--- антицепь размера $k + 1$~--- противоречие.
		
		\textit{Почему $X^+ \cap X^- = A$:} если $x \in X^+ \cap X^-$, то $x \geq a$ и $x \leq b$ для некоторых $a, b \in A$. Тогда $a \leq x \leq b$, и раз $A$~--- антицепь, $a = b$, откуда $x = a \in A$.
		
		Так как $A$ не состоит из всех максимальных элементов $X$, $|X^+| < |X|$ (есть элемент $X$, не лежащий в $X^+$). Аналогично $|X^-| < |X|$.
		
		Максимальный размер антицепи в $X^+$: не превосходит $k$ (любая антицепь $X^+$ является антицепью $X$); с другой стороны, $A \subseteq X^+$ является антицепью размера $k$.
		
		По индукции $X^+$ разбивается на $k$ цепей $C_1^+, \dots, C_k^+$. Каждый элемент $A$ лежит в $X^+$, и в каждой цепи $C_i^+$ есть хотя бы один элемент $A$~--- не более одного (антицепь). Значит, каждая $C_i^+$ содержит ровно по одному элементу $A$.
		
		Аналогично $X^-$ разбивается на $k$ цепей $C_1^-, \dots, C_k^-$, каждая содержит ровно один элемент $A$.
		
		Нумеруем так, чтобы $C_i^+$ и $C_i^-$ содержали один и тот же элемент $a_i \in A$. Тогда $C_i = C_i^- \cup C_i^+$~--- цепь: любые $x \in C_i^-$ и $y \in C_i^+$ удовлетворяют $x \leq a_i \leq y$.
		
		Цепи $C_1, \dots, C_k$ покрывают $X = X^- \cup X^+$.
		
		\textbf{Случай 2:} каждая максимальная антицепь размера $k$ состоит только из максимальных элементов или только из минимальных.
		
		Возьмём максимальную цепь $C$ с минимальным элементом $m$ и максимальным $M$. Рассмотрим $X' = X \setminus C$.
		
		Утверждение: максимальный размер антицепи в $X'$ равен $k - 1$.
		
		\textit{Не превышает $k$:} очевидно.
		
		\textit{Если бы равнялся $k$:} пусть $A'$~--- антицепь в $X'$ размера $k$. Она является антицепью в $X$, следовательно, по условию случая 2, состоит только из максимальных (или только из минимальных) элементов $X$. Если $A'$ состоит из максимальных, добавим $m$ (минимальный элемент $C$~--- он не является максимальным при $|C| > 1$, а если $|C| = 1$, то $m = M$, и мы добавляем этот элемент к другой антицепи). На самом деле, $m$ несравним с элементами $A'$: если бы $m \leq a'$ для $a' \in A'$, то $a'$ не максимален, противоречие. Итак, $A' \cup \{m\}$~--- антицепь размера $k + 1$~--- противоречие.
		
		Значит в $X'$ максимальная антицепь имеет размер $\leq k - 1$. Но $k - 1$ достигается (иначе бы $X' = X \setminus C$ покрывалось $k - 2$ цепями, и вместе с $C$ мы бы покрыли $X$ ровно $k - 1$ цепью, противореча нижней оценке).
		
		По индукции $X'$ разбивается на $k - 1$ цепей. Вместе с $C$ получаем $k$ цепей, покрывающих $X$.
	\end{proof}
	
	\begin{theorem}[Мирский]
		Пусть $(X, \leq)$~--- конечное частично упорядоченное множество, и максимальный размер цепи равен $\ell$. Тогда $X$ можно разбить на $\ell$ антицепей, и не менее.
	\end{theorem}
	
	\begin{proof}
		Не менее $\ell$ антицепей необходимо: каждая антицепь содержит не более одного элемента любой цепи.
		
		Определим \textbf{ранг} элемента $x$:
		\[r(x) = \text{(длина максимальной цепи с верхним концом в } x) - 1.\]
		
		Иными словами, $r(x) = 0$ если $x$ минимален (нет элементов строго меньше), $r(x) = \max_{y < x} r(y) + 1$ в общем случае.
		
		Пусть $A_i = \{x \in X : r(x) = i\}$ для $i = 0, 1, \dots, \ell - 1$.
		
		\textbf{$A_i$~--- антицепь:} если $x, y \in A_i$ и $x < y$, то максимальная цепь, оканчивающаяся в $x$ (длины $i + 1$), вместе с $y$ даёт цепь длины $i + 2$. Тогда $r(y) \geq i + 1$, противоречие с $y \in A_i$.
		
		\textbf{Покрытие:} каждый $x \in X$ имеет ранг от $0$ до $\ell - 1$, так что $X = A_0 \sqcup A_1 \sqcup \dots \sqcup A_{\ell - 1}$.
		
		\textbf{Ровно $\ell$ антицепей:} так как максимальная цепь имеет длину $\ell$, её элементы $x_0 < x_1 < \dots < x_{\ell-1}$ имеют ранги $0, 1, \dots, \ell - 1$ соответственно (ранг возрастает вдоль цепи), значит все $A_i$ непусты.
	\end{proof}
	
	\begin{corollary}[Теорема Эрдёша--Секереша]
		В последовательности из $n > (r-1)(s-1)$ различных чисел найдётся возрастающая подпоследовательность длины $r$ или убывающая подпоследовательность длины $s$.
	\end{corollary}
	
	\begin{proof}
		Рассмотрим частичный порядок на позициях последовательности $a_1, \dots, a_n$: $i \leq j \Leftrightarrow i \leq j$ (как индексы) и $a_i \leq a_j$ (как числа).
		
		Цепи в этом порядке~--- в точности возрастающие подпоследовательности, антицепи~--- убывающие подпоследовательности (если $i < j$ и $a_i > a_j$, то $i$ и $j$ несравнимы).
		
		Если нет возрастающей подпоследовательности длины $r$, то максимальная длина цепи $\leq r - 1$.
		
		По теореме Мирского, $X$ разбивается на $\leq r - 1$ антицепей.
		
		По принципу Дирихле, одна из них содержит $\geq \lceil n/(r-1) \rceil > (r-1)(s-1)/(r-1) = s - 1$ элементов, то есть $\geq s$ элементов~--- это убывающая подпоследовательность длины $\geq s$.
	\end{proof}
	
	\newpage
	
	\section{Мощности множеств. Счётные множества и их свойства. Несуществование мощности между конечными и счётными. Множества мощности континуум, свойства. Несуществование наибольшей мощности. Континуум как мощность множества подмножеств счётного множества. Теорема Кантора-Бернштейна.}
	
	\subsection{Мощности множеств}
	
	\begin{definition}
		Множества $A$ и $B$ \textbf{равномощны} ($|A| = |B|$ или $A \sim B$), если существует биекция $f: A \to B$.
	\end{definition}
	
	\begin{definition}
		$|A| \leqslant |B|$, если существует инъекция $f: A \to B$. $|A| < |B|$, если $|A| \leqslant |B|$ и $|A| \neq |B|$.
	\end{definition}
	
	\begin{note}
		Равномощность~--- отношение эквивалентности (рефлексивность: $\text{id}$; симметричность: обратная биекция; транзитивность: композиция биекций).
	\end{note}
	
	\subsection{Счётные множества}
	
	\begin{definition}
		Множество называется \textbf{счётным}, если оно равномощно $\mathbb{N}$.
		Множество называется \textbf{не более чем счётным}, если оно конечно или счётно.
	\end{definition}
	
	\begin{theorem}[Свойства счётных множеств]
		\begin{enumerate}
			\item Подмножество счётного множества конечно или счётно.
			\item Объединение счётного числа счётных множеств счётно.
			\item $\mathbb{Z}$ счётно.
			\item $\mathbb{Q}$ счётно.
			\item $\mathbb{N} \times \mathbb{N}$ счётно.
		\end{enumerate}
	\end{theorem}
	
	\begin{proof}
		1. Пусть $A \subseteq B$, $B = \{b_1, b_2, \dots\}$ счётно. Если $A$ бесконечно, перенумеруем его элементы в порядке их появления в $B$: $a_1 = b_{n_1}, a_2 = b_{n_2}, \dots$, где $n_1 < n_2 < \dots$ Отображение $k \mapsto a_k$ является биекцией $\mathbb{N} \to A$.
		
		2. Пусть $A_1, A_2, \dots$~--- счётные, $A_i = \{a_{i1}, a_{i2}, \dots\}$. Запишем элементы в таблицу:
		\[a_{11}, a_{12}, a_{13}, \dots\]
		\[a_{21}, a_{22}, a_{23}, \dots\]
		\[\vdots\]
		
		Обходим по диагоналям: $a_{11}, a_{12}, a_{21}, a_{13}, a_{22}, a_{31}, \dots$ Это даёт сюръекцию $\mathbb{N} \to \bigcup A_i$ (некоторые элементы могут повторяться, если множества не дизъюнктны; выкидываем повторения). По пункту 1, бесконечное подмножество счётного счётно.
		
		3. Биекция $\mathbb{N} \to \mathbb{Z}$: $n \mapsto \begin{cases} n/2, & n \text{ чётно} \\ -(n+1)/2, & n \text{ нечётно} \end{cases}$, то есть $0, 1, -1, 2, -2, \dots$
		
		4. $\mathbb{Q}^+ = \{p/q : p, q \in \mathbb{N}, \gcd(p,q) = 1\} \subseteq \mathbb{N} \times \mathbb{N}$~--- счётно (как подмножество счётного). $\mathbb{Q} = \mathbb{Q}^+ \cup \{0\} \cup (-\mathbb{Q}^+)$~--- объединение трёх счётных (или конечных) множеств~--- счётно.
		
		5. Явная биекция: пара $(m, n)$ отображается в $\frac{(m+n-2)(m+n-1)}{2} + m$ (порядковый номер при обходе по диагоналям). Или через функцию Кантора: $(m, n) \mapsto \frac{(m+n)(m+n+1)}{2} + n$.
	\end{proof}
	
	\subsection{Несуществование мощности между конечными и счётными}
	
	\begin{theorem}
		Любое бесконечное множество содержит счётное подмножество.
	\end{theorem}
	
	\begin{proof}
		Пусть $A$ бесконечно. Выберем $a_1 \in A$ (возможно, т.к. $A \neq \varnothing$). Затем $a_2 \in A \setminus \{a_1\}$ (возможно, т.к. $A$ бесконечно, значит $A \neq \{a_1\}$). Продолжаем: на $n$-м шаге $a_n \in A \setminus \{a_1, \dots, a_{n-1}\}$ (возможно, т.к. $A$ бесконечно, значит $A \neq \{a_1, \dots, a_{n-1}\}$). Получаем счётное множество $\{a_1, a_2, \dots\} \subseteq A$, где все $a_i$ различны.
	\end{proof}
	
	\begin{corollary}
		Не существует мощности, строго большей всех конечных и строго меньшей $|\mathbb{N}|$: любое бесконечное множество имеет мощность $\geqslant |\mathbb{N}|$.
	\end{corollary}
	
	\begin{proof}
		Из теоремы: любое бесконечное $A$ содержит счётное $B \subseteq A$. Значит существует инъекция $\mathbb{N} \to B \hookrightarrow A$, то есть $|\mathbb{N}| \leqslant |A|$.
	\end{proof}
	
	\subsection{Континуум}
	
	\begin{definition}
		Мощность $\mathfrak{c} = |\mathbb{R}|$ называется \textbf{континуумом}.
	\end{definition}
	
	\begin{theorem}
		$|(0, 1)| = |\mathbb{R}|$.
	\end{theorem}
	
	\begin{proof}
		Биекция $(0, 1) \to \mathbb{R}$: $f(x) = \tan\!\left(\pi x - \frac{\pi}{2}\right)$.
		
		Это строго возрастающая непрерывная функция, $f(0^+) = -\infty$, $f(1^-) = +\infty$.
	\end{proof}
	
	\begin{theorem}[Кантор]
		$\mathbb{R}$ несчётно, то есть $|\mathbb{R}| > |\mathbb{N}|$.
	\end{theorem}
	
	\begin{proof}[Диагональный аргумент]
		Предположим, что $(0, 1)$ счётно: $(0, 1) = \{r_1, r_2, r_3, \dots\}$, где каждое $r_i$ записано в виде десятичной дроби $r_i = 0.d_{i1}d_{i2}d_{i3}\dots$ (цифры $d_{ij} \in \{0, 1, \dots, 9\}$).
		
		Построим число $r = 0.d_1d_2d_3\dots$, где:
		\[d_i = \begin{cases} 1, & d_{ii} \neq 1 \\ 2, & d_{ii} = 1 \end{cases}\]
		
		(Выбираем $d_i \notin \{0, 9\}$, чтобы избежать неоднозначности десятичной записи.)
		
		Тогда $r \in (0, 1)$, но $r \neq r_i$ для каждого $i$~--- ведь $r$ и $r_i$ отличаются в $i$-й десятичной цифре: $d_i \neq d_{ii}$.
		
		Это противоречит тому, что $(0, 1) = \{r_1, r_2, \dots\}$.
	\end{proof}
	
	\begin{theorem}[Свойства континуума]
		\begin{enumerate}
			\item $|[0,1]| = |(0,1)| = |\mathbb{R}|$
			\item $|\mathbb{R}^n| = |\mathbb{R}|$ для любого $n \geqslant 1$
			\item $|\mathbb{R}| = |2^{\mathbb{N}}|$ (континуум~--- это мощность множества всех подмножеств натуральных чисел)
		\end{enumerate}
	\end{theorem}
	
	\begin{proof}
		1. $(0,1) \subseteq [0,1] \subseteq \mathbb{R}$. Имеем инъекции $(0,1) \hookrightarrow [0,1] \hookrightarrow \mathbb{R}$, то есть $|(0,1)| \leqslant |[0,1]| \leqslant |\mathbb{R}|$. С другой стороны, $|(0,1)| = |\mathbb{R}|$. По теореме Кантора--Бернштейна $|[0,1]| = |\mathbb{R}|$.
		
		2. Достаточно для $n = 2$ (общий случай~--- индукция). Покажем $|(0,1)^2| = |(0,1)|$.
		
		Инъекция $(0,1) \to (0,1)^2$: $x \mapsto (x, 1/2)$.
		
		Инъекция $(0,1)^2 \to (0,1)$: запишем $x = 0.a_1a_2a_3\dots$ и $y = 0.b_1b_2b_3\dots$ без девяток в хвосте (устраняем неоднозначность). Сопоставим $(x, y) \mapsto 0.a_1b_1a_2b_2a_3b_3\dots$~--- это инъекция.
		
		По теореме Кантора--Бернштейна, $|(0,1)^2| = |(0,1)|$.
		
		3. Покажем $|\mathbb{R}| = |2^\mathbb{N}|$.
		
		\textit{Инъекция $2^\mathbb{N} \to [0,1]$:} подмножеству $A \subseteq \mathbb{N}$ сопоставим $\sum_{n \in A} 2 \cdot 3^{-n}$ (число в тернарной записи с цифрами 0 и 2). Это инъекция (числа Кантора различны).
		
		\textit{Инъекция $[0,1] \to 2^\mathbb{N}$:} числу $x \in [0,1]$ в двоичной записи $x = 0.b_1b_2b_3\dots$ (при неоднозначности берём запись с нулями в хвосте) сопоставим $\{n \in \mathbb{N} : b_n = 1\} \subseteq \mathbb{N}$.
		
		По теореме Кантора--Бернштейна, $|[0,1]| = |2^\mathbb{N}|$, и по пункту 1, $|\mathbb{R}| = |2^\mathbb{N}|$.
	\end{proof}
	
	\subsection{Несуществование наибольшей мощности}
	
	\begin{theorem}[Кантор]
		Для любого множества $A$: $|A| < |2^A|$ (строгое неравенство).
	\end{theorem}
	
	\begin{proof}
		\textit{$|A| \leqslant |2^A|$:} инъекция $A \to 2^A$: $a \mapsto \{a\}$.
		
		\textit{$|A| \neq |2^A|$:} предположим, существует биекция $f: A \to 2^A$. Рассмотрим множество (диагональное построение):
		\[B = \{a \in A : a \notin f(a)\}.\]
		
		$B \subseteq A$, значит $B \in 2^A$. Так как $f$ биекция, существует $b \in A$: $f(b) = B$.
		
		Вопрос: $b \in B$?
		\begin{itemize}
			\item Если $b \in B$, то по определению $B$: $b \notin f(b) = B$. Противоречие.
			\item Если $b \notin B$, то $b \notin f(b)$, значит по определению $B$: $b \in B$. Противоречие.
		\end{itemize}
		
		В обоих случаях противоречие. Значит, биекции $f: A \to 2^A$ не существует.
	\end{proof}
	
	\begin{corollary}
		Существует бесконечная строгая цепочка мощностей:
		\[|\mathbb{N}| < |2^{\mathbb{N}}| < |2^{2^\mathbb{N}}| < \dots\]
	\end{corollary}
	
	\subsection{Теорема Кантора--Бернштейна}
	
	\begin{theorem}[Кантора--Бернштейна]
		Если $|A| \leqslant |B|$ и $|B| \leqslant |A|$, то $|A| = |B|$.
	\end{theorem}
	
	\begin{proof}
		Пусть $f: A \to B$ и $g: B \to A$~--- инъекции. Построим биекцию $h: A \to B$.
		
		\textit{Идея:} разобьём $A$ на две части: на первой $h = f$, на второй $h = g^{-1}$.
		
		Определим последовательность множеств:
		\[A_0 = A \setminus g(B), \quad A_{n+1} = g(f(A_n)).\]
		
		$A_0$~--- это «элементы $A$, которые не являются образом ничего из $B$» (они «начинают новую нить»). $A_{n+1}$~--- образы образов.
		
		Пусть $A^* = \bigcup_{n=0}^{\infty} A_n$.
		
		Определим $h: A \to B$:
		\[h(a) = \begin{cases} f(a), & \text{если } a \in A^* \\ g^{-1}(a), & \text{если } a \notin A^* \end{cases}\]
		
		\textit{Корректность $g^{-1}(a)$ при $a \notin A^*$:} если $a \notin A^*$, то в частности $a \notin A_0 = A \setminus g(B)$, значит $a \in g(B)$, то есть существует (единственный, т.к. $g$ инъекция) элемент $b \in B$ с $g(b) = a$. Полагаем $g^{-1}(a) = b$.
		
		\textit{Инъективность $h$:}
		\begin{itemize}
			\item На $A^*$ функция $h = f$ инъективна.
			\item На $A \setminus A^*$ функция $h = g^{-1}$ инъективна (как ограничение инъекции).
			\item Пусть $a \in A^*$, $a' \notin A^*$ и $h(a) = h(a')$. Тогда $f(a) = g^{-1}(a')$, то есть $g(f(a)) = a'$. Но $a \in A^* = \bigcup A_n$, значит $a \in A_n$ для некоторого $n$, откуда $a' = g(f(a)) \in g(f(A_n)) = A_{n+1} \subseteq A^*$~--- противоречие с $a' \notin A^*$.
		\end{itemize}
		
		\textit{Сюръективность $h$:} для $b \in B$:
		\begin{itemize}
			\item Если $g(b) \in A^*$, то $g(b) \in A_n$ для некоторого $n$. Так как $g(b) \in g(B)$, а $A_0 = A \setminus g(B)$, имеем $g(b) \notin A_0$, значит $n \geqslant 1$. Тогда $g(b) \in A_n = g(f(A_{n-1}))$, откуда $g(b) = g(f(a))$ для некоторого $a \in A_{n-1} \subseteq A^*$. По инъективности $g$: $b = f(a)$. Тогда $h(a) = f(a) = b$.
			\item Если $g(b) \notin A^*$, то $h(g(b)) = g^{-1}(g(b)) = b$.
		\end{itemize}
		
		Таким образом, $h$~--- биекция.
	\end{proof}
	
	\newpage
	
	\section{Дискретное вероятностное пространство. Парадокс Бертрана. Вероятности событий, независимые события. Формулы полной вероятности и Байеса. Случайные величины и их функции распределения. Математическое ожидание и его свойства. Дисперсия и её свойства. Классические распределения: биномиальное, геометрическое, гипергеометрическое. Распределение Пуассона как предел биномиального. Математическое ожидание и дисперсия классических распределений. Неравенство Чебышёва и закон больших чисел.}
	
	\subsection{Дискретное вероятностное пространство}
	
	\begin{definition}
		\textbf{Дискретное вероятностное пространство}~--- пара $(\Omega, P)$, где $\Omega$~--- конечное или счётное множество элементарных исходов, $P: \Omega \to [0, 1]$ такова, что $\sum_{\omega \in \Omega} P(\omega) = 1$.
	\end{definition}
	
	\begin{definition}
		\textbf{Событие}~--- подмножество $A \subseteq \Omega$. Вероятность события: $P(A) = \sum_{\omega \in A} P(\omega)$.
	\end{definition}
	
	\begin{theorem}[Свойства вероятности]
		\begin{enumerate}
			\item $P(\varnothing) = 0$, $P(\Omega) = 1$
			\item $P(A^c) = 1 - P(A)$
			\item $P(A \cup B) = P(A) + P(B) - P(A \cap B)$ (формула включений-исключений)
			\item Если $A \cap B = \varnothing$, то $P(A \cup B) = P(A) + P(B)$ (аддитивность)
		\end{enumerate}
	\end{theorem}
	
	\begin{proof}
		1. $P(\varnothing) = \sum_{\omega \in \varnothing} P(\omega) = 0$; $P(\Omega) = \sum_{\omega \in \Omega} P(\omega) = 1$ по определению.
		
		2. $1 = P(\Omega) = P(A \cup A^c) = P(A) + P(A^c)$ (т.к. $A \cap A^c = \varnothing$).
		
		3. $A \cup B = (A \setminus B) \sqcup (A \cap B) \sqcup (B \setminus A)$, значит $P(A \cup B) = P(A \setminus B) + P(A \cap B) + P(B \setminus A)$. Из $A = (A \setminus B) \sqcup (A \cap B)$: $P(A \setminus B) = P(A) - P(A \cap B)$. Аналогично $P(B \setminus A) = P(B) - P(A \cap B)$.
	\end{proof}
	
	\subsection{Парадокс Бертрана}
	
	\textbf{Задача:} Какова вероятность того, что случайная хорда окружности длиннее стороны вписанного равностороннего треугольника?
	
	\textbf{Три модели:}
	\begin{enumerate}
		\item \textbf{Случайные концы:} Фиксируем один конец хорды, второй выбираем равномерно на окружности. Хорда длиннее стороны треугольника, если второй конец попадает в дугу длиной $1/3$ окружности. $P = 1/3$.
		\item \textbf{Случайная середина на диаметре:} Выбираем диаметр, затем середину хорды равномерно на диаметре. Хорда длиннее, если середина ближе к центру, чем $r/2$, то есть на отрезке длиной $r$ из $2r$. $P = 1/2$.
		\item \textbf{Случайная середина в круге:} Середина хорды выбирается равномерно в круге. Хорда длиннее, если середина попадает в круг радиуса $r/2$, то есть площадью $\pi(r/2)^2 = \pi r^2/4$. $P = 1/4$.
	\end{enumerate}
	
	\textbf{Мораль:} Вероятностная модель должна быть чётко определена. «Случайная хорда» не имеет однозначного смысла без указания, какое именно вероятностное пространство используется.
	
	\subsection{Независимые события}
	
	\begin{definition}
		События $A$ и $B$ \textbf{независимы}, если $P(A \cap B) = P(A) \cdot P(B)$.
	\end{definition}
	
	\begin{note}
		Интуиция: $A$ и $B$ независимы, если знание о том, что $B$ произошло, не меняет вероятность $A$, то есть $P(A \mid B) = P(A)$ (при $P(B) > 0$).
	\end{note}
	
	\begin{definition}
		События $A_1, \dots, A_n$ \textbf{независимы в совокупности}, если для любого $I \subseteq \{1, \dots, n\}$:
		\[P\!\left(\bigcap_{i \in I} A_i\right) = \prod_{i \in I} P(A_i).\]
	\end{definition}
	
	\begin{definition}
		\textbf{Условная вероятность:} $P(A \mid B) = \dfrac{P(A \cap B)}{P(B)}$ при $P(B) > 0$.
	\end{definition}
	
	\begin{theorem}[Формула умножения]
		$P(A_1 \cap A_2 \cap \dots \cap A_n) = P(A_1) \cdot P(A_2 \mid A_1) \cdot P(A_3 \mid A_1 \cap A_2) \cdots$
	\end{theorem}
	
	\begin{proof}
		Из определения условной вероятности: $P(A \cap B) = P(A) \cdot P(B \mid A)$. Применяем последовательно.
	\end{proof}
	
	\subsection{Формулы полной вероятности и Байеса}
	
	\begin{theorem}[Формула полной вероятности]
		Пусть $B_1, \dots, B_n$~--- разбиение $\Omega$ ($B_i \cap B_j = \varnothing$ при $i \neq j$, $\bigcup B_i = \Omega$), $P(B_i) > 0$. Тогда:
		\[P(A) = \sum_{i=1}^{n} P(A \mid B_i) \cdot P(B_i).\]
	\end{theorem}
	
	\begin{proof}
		$A = A \cap \Omega = A \cap \bigsqcup_{i} B_i = \bigsqcup_{i} (A \cap B_i)$.
		
		Из аддитивности вероятности:
		\[P(A) = \sum_{i=1}^{n} P(A \cap B_i) = \sum_{i=1}^{n} P(A \mid B_i) \cdot P(B_i).\]
	\end{proof}
	
	\begin{theorem}[Формула Байеса]
		При тех же условиях и при $P(A) > 0$:
		\[P(B_k \mid A) = \frac{P(A \mid B_k) \cdot P(B_k)}{\displaystyle\sum_{i=1}^{n} P(A \mid B_i) \cdot P(B_i)}.\]
	\end{theorem}
	
	\begin{proof}
		По определению условной вероятности и формуле умножения:
		\[P(B_k \mid A) = \frac{P(A \cap B_k)}{P(A)} = \frac{P(A \mid B_k) \cdot P(B_k)}{P(A)}.\]
		Знаменатель раскрывается по формуле полной вероятности.
	\end{proof}
	
	\begin{note}
		$P(B_k)$ называется \textbf{априорной} вероятностью, $P(B_k \mid A)$~--- \textbf{апостериорной} (после наблюдения события $A$). Байесовский подход: «обновление» вероятности гипотезы $B_k$ после получения данных $A$.
	\end{note}
	
	\subsection{Случайные величины}
	
	\begin{definition}
		\textbf{Случайная величина}~--- функция $X: \Omega \to \mathbb{R}$.
	\end{definition}
	
	\begin{definition}
		\textbf{Функция распределения:} $F_X(x) = P(X \leqslant x) = P(\{\omega \in \Omega : X(\omega) \leqslant x\})$.
	\end{definition}
	
	\begin{note}
		Функция распределения:
		\begin{itemize}
			\item монотонно не убывает,
			\item $\lim_{x \to -\infty} F_X(x) = 0$, $\lim_{x \to +\infty} F_X(x) = 1$,
			\item непрерывна справа.
		\end{itemize}
		В дискретном случае $F_X$ кусочно постоянна с разрывами в точках $x$, где $P(X = x) > 0$.
	\end{note}
	
	\begin{definition}
		Функция вероятности (закон распределения): $p_X(x) = P(X = x)$.
	\end{definition}
	
	\subsection{Математическое ожидание}
	
	\begin{definition}
		\textbf{Математическое ожидание} (среднее):
		\[\mathbb{E}[X] = \sum_{\omega \in \Omega} X(\omega) P(\omega) = \sum_{x} x \cdot P(X = x),\]
		при условии абсолютной сходимости ряда.
	\end{definition}
	
	\begin{theorem}[Свойства математического ожидания]
		\begin{enumerate}
			\item \textbf{Линейность:} $\mathbb{E}[aX + bY] = a\mathbb{E}[X] + b\mathbb{E}[Y]$
			\item \textbf{Монотонность:} если $X \geqslant 0$, то $\mathbb{E}[X] \geqslant 0$; если $X \leqslant Y$, то $\mathbb{E}[X] \leqslant \mathbb{E}[Y]$
			\item \textbf{Независимость:} если $X$ и $Y$ независимы, то $\mathbb{E}[XY] = \mathbb{E}[X] \cdot \mathbb{E}[Y]$
		\end{enumerate}
	\end{theorem}
	
	\begin{proof}
		1. $\mathbb{E}[aX + bY] = \sum_\omega (aX(\omega) + bY(\omega)) P(\omega) = a \sum_\omega X(\omega) P(\omega) + b \sum_\omega Y(\omega) P(\omega) = a\mathbb{E}[X] + b\mathbb{E}[Y]$.
		
		2. $\mathbb{E}[X] = \sum_\omega X(\omega) P(\omega) \geqslant 0$ при $X(\omega) \geqslant 0$.
		
		3. $X$ и $Y$ независимы означает $P(X = x, Y = y) = P(X = x) \cdot P(Y = y)$. Тогда:
		\[\mathbb{E}[XY] = \sum_{x,y} xy \cdot P(X=x, Y=y) = \sum_{x,y} xy \cdot P(X=x) P(Y=y) = \Bigl(\sum_x x P(X=x)\Bigr) \Bigl(\sum_y y P(Y=y)\Bigr) = \mathbb{E}[X]\mathbb{E}[Y].\]
	\end{proof}
	
	\begin{theorem}[Формула замены переменной]
		$\mathbb{E}[g(X)] = \sum_x g(x) P(X = x)$.
	\end{theorem}
	
	\begin{proof}
		$\mathbb{E}[g(X)] = \sum_\omega g(X(\omega)) P(\omega) = \sum_x g(x) \sum_{\omega: X(\omega) = x} P(\omega) = \sum_x g(x) P(X = x)$.
	\end{proof}
	
	\subsection{Дисперсия}
	
	\begin{definition}
		\textbf{Дисперсия:} $\mathbb{D}[X] = \mathbb{E}[(X - \mathbb{E}[X])^2]$.
	\end{definition}
	
	\begin{theorem}[Вычислительная формула]
		$\mathbb{D}[X] = \mathbb{E}[X^2] - (\mathbb{E}[X])^2$.
	\end{theorem}
	
	\begin{proof}
		Пусть $\mu = \mathbb{E}[X]$. Тогда:
		\begin{align*}
			\mathbb{D}[X] &= \mathbb{E}[(X - \mu)^2] = \mathbb{E}[X^2 - 2\mu X + \mu^2] \\
			&= \mathbb{E}[X^2] - 2\mu \mathbb{E}[X] + \mu^2 = \mathbb{E}[X^2] - 2\mu^2 + \mu^2 = \mathbb{E}[X^2] - \mu^2.
		\end{align*}
	\end{proof}
	
	\begin{theorem}[Свойства дисперсии]
		\begin{enumerate}
			\item $\mathbb{D}[X] \geqslant 0$; $\mathbb{D}[X] = 0 \Leftrightarrow X = \mathbb{E}[X]$ почти наверное
			\item $\mathbb{D}[aX + b] = a^2 \mathbb{D}[X]$
			\item Если $X, Y$ независимы: $\mathbb{D}[X + Y] = \mathbb{D}[X] + \mathbb{D}[Y]$
		\end{enumerate}
	\end{theorem}
	
	\begin{proof}
		1. $\mathbb{D}[X] = \mathbb{E}[(X - \mathbb{E}[X])^2] \geqslant 0$, т.к. квадрат неотрицателен. Равенство нулю означает $(X - \mu)^2 = 0$ везде, то есть $X = \mu$.
		
		2. $\mathbb{E}[aX+b] = a\mu + b$, поэтому:
		\[\mathbb{D}[aX+b] = \mathbb{E}[(aX+b - (a\mu+b))^2] = \mathbb{E}[a^2(X-\mu)^2] = a^2 \mathbb{D}[X].\]
		
		3. При независимости $\mathbb{E}[XY] = \mathbb{E}[X]\mathbb{E}[Y]$. Пусть $\mu_X = \mathbb{E}[X]$, $\mu_Y = \mathbb{E}[Y]$:
		\begin{align*}
			\mathbb{D}[X+Y] &= \mathbb{E}[(X+Y)^2] - (\mu_X + \mu_Y)^2 \\
			&= \mathbb{E}[X^2] + 2\mathbb{E}[XY] + \mathbb{E}[Y^2] - \mu_X^2 - 2\mu_X\mu_Y - \mu_Y^2 \\
			&= (\mathbb{E}[X^2] - \mu_X^2) + 2(\mathbb{E}[XY] - \mu_X\mu_Y) + (\mathbb{E}[Y^2] - \mu_Y^2) \\
			&= \mathbb{D}[X] + 0 + \mathbb{D}[Y].
		\end{align*}
	\end{proof}
	
	\begin{definition}
		\textbf{Ковариация:} $\text{Cov}(X, Y) = \mathbb{E}[(X - \mathbb{E}[X])(Y - \mathbb{E}[Y])] = \mathbb{E}[XY] - \mathbb{E}[X]\mathbb{E}[Y]$.
		
		При независимости $X$ и $Y$: $\text{Cov}(X, Y) = 0$ (но не наоборот).
	\end{definition}
	
	\subsection{Классические распределения}
	
	\textbf{Биномиальное распределение} $X \sim \text{Bin}(n, p)$:
	\[P(X = k) = \binom{n}{k} p^k (1-p)^{n-k}, \quad k = 0, 1, \dots, n.\]
	
	\textit{Интерпретация:} число успехов в $n$ независимых испытаниях Бернулли с вероятностью успеха $p$.
	
	\textit{Проверка нормировки:} $\sum_{k=0}^{n} \binom{n}{k} p^k (1-p)^{n-k} = (p + (1-p))^n = 1$.
	
	\begin{theorem}
		$\mathbb{E}[X] = np$, $\mathbb{D}[X] = np(1-p)$.
	\end{theorem}
	
	\begin{proof}
		Представим $X = X_1 + \dots + X_n$, где $X_i$~--- индикатор успеха в $i$-м испытании:
		\[X_i = \begin{cases} 1, & \text{успех в } i\text{-м испытании} \\ 0, & \text{неудача} \end{cases}\]
		
		$\mathbb{E}[X_i] = 1 \cdot p + 0 \cdot (1-p) = p$.
		
		$\mathbb{E}[X_i^2] = 1^2 \cdot p + 0^2 \cdot (1-p) = p$, откуда $\mathbb{D}[X_i] = p - p^2 = p(1-p)$.
		
		По линейности: $\mathbb{E}[X] = \sum_{i=1}^n \mathbb{E}[X_i] = np$.
		
		Так как $X_1, \dots, X_n$ независимы: $\mathbb{D}[X] = \sum_{i=1}^n \mathbb{D}[X_i] = np(1-p)$.
	\end{proof}
	
	\textbf{Геометрическое распределение} $X \sim \text{Geom}(p)$:
	\[P(X = k) = (1-p)^{k-1} p, \quad k = 1, 2, 3, \dots\]
	
	\textit{Интерпретация:} номер первого успеха в последовательности независимых испытаний Бернулли.
	
	\textit{Проверка нормировки:} $\sum_{k=1}^{\infty} (1-p)^{k-1} p = p \cdot \dfrac{1}{1-(1-p)} = 1$.
	
	\begin{theorem}
		$\mathbb{E}[X] = \dfrac{1}{p}$, $\mathbb{D}[X] = \dfrac{1-p}{p^2}$.
	\end{theorem}
	
	\begin{proof}
		Обозначим $q = 1 - p$.
		
		\textit{Математическое ожидание:}
		\[\mathbb{E}[X] = \sum_{k=1}^{\infty} k q^{k-1} p = p \sum_{k=1}^{\infty} k q^{k-1} = p \cdot \frac{d}{dq}\sum_{k=1}^{\infty} q^k = p \cdot \frac{d}{dq} \frac{q}{1-q} = p \cdot \frac{1}{(1-q)^2} = \frac{p}{p^2} = \frac{1}{p}.\]
		
		\textit{Второй момент:}
		\[\mathbb{E}[X^2] = \sum_{k=1}^{\infty} k^2 q^{k-1} p.\]
		
		Используем $k^2 = k(k-1) + k$:
		\[\mathbb{E}[X^2] = p\sum_{k=1}^{\infty} k(k-1) q^{k-1} + p\sum_{k=1}^{\infty} k q^{k-1} = p \sum_{k=2}^{\infty} k(k-1) q^{k-1} + \frac{1}{p}.\]
		
		$\sum_{k=2}^{\infty} k(k-1) q^{k-1} = \frac{d}{dq}\sum_{k=2}^{\infty} k q^k = \frac{d}{dq}\left(q \cdot \frac{d}{dq}\sum_{k=1}^{\infty} q^k\right) = \frac{d}{dq}\frac{q}{(1-q)^2} = \frac{1+q}{(1-q)^3}$.
		
		Поэтому $\mathbb{E}[X^2] = p \cdot \frac{1+q}{p^3} + \frac{1}{p} = \frac{1+q}{p^2} + \frac{1}{p} = \frac{1+q+p}{p^2} = \frac{2-p}{p^2}$.
		
		\[\mathbb{D}[X] = \mathbb{E}[X^2] - (\mathbb{E}[X])^2 = \frac{2-p}{p^2} - \frac{1}{p^2} = \frac{1-p}{p^2} = \frac{q}{p^2}.\]
	\end{proof}
	
	\textbf{Гипергеометрическое распределение:}
	
	Из $N$ объектов $K$ отмечены. Выбираем $n$ без возвращения. $X$~--- число отмеченных.
	\[P(X = k) = \frac{\binom{K}{k} \binom{N-K}{n-k}}{\binom{N}{n}}, \quad \max(0, n-(N-K)) \leqslant k \leqslant \min(n, K).\]
	
	\textit{Комбинаторный смысл:} $\binom{K}{k}$~--- число способов выбрать $k$ отмеченных, $\binom{N-K}{n-k}$~--- число способов выбрать $n-k$ неотмеченных.
	
	\begin{theorem}
		$\mathbb{E}[X] = \dfrac{nK}{N}$.
	\end{theorem}
	
	\begin{proof}
		Представим $X = X_1 + \dots + X_n$, где $X_i$~--- индикатор того, что $i$-й выбранный объект отмечен.
		
		$\mathbb{E}[X_i] = P(X_i = 1) = \dfrac{K}{N}$ (вероятность того, что конкретный объект отмечен~--- по симметрии).
		
		По линейности: $\mathbb{E}[X] = n \cdot \dfrac{K}{N} = \dfrac{nK}{N}$.
	\end{proof}
	
	\begin{note}
		$\mathbb{D}[X] = \dfrac{nK(N-K)(N-n)}{N^2(N-1)}$~--- вычисление аналогично, через $\mathbb{E}[X_i X_j]$.
	\end{note}
	
	\subsection{Распределение Пуассона}
	
	\begin{definition}
		$X \sim \text{Pois}(\lambda)$: $P(X = k) = \dfrac{\lambda^k e^{-\lambda}}{k!}$, $k = 0, 1, 2, \dots$
	\end{definition}
	
	\textit{Проверка нормировки:} $\sum_{k=0}^{\infty} \dfrac{\lambda^k e^{-\lambda}}{k!} = e^{-\lambda} e^{\lambda} = 1$.
	
	\begin{theorem}
		При $n \to \infty$, $p = p_n \to 0$, $np_n \to \lambda$: $\text{Bin}(n, p_n) \xrightarrow{d} \text{Pois}(\lambda)$, то есть для каждого $k = 0, 1, 2, \dots$:
		\[\binom{n}{k} p_n^k (1-p_n)^{n-k} \to \frac{\lambda^k e^{-\lambda}}{k!}.\]
	\end{theorem}
	
	\begin{proof}
		Фиксируем $k$. Пусть $\lambda_n = np_n \to \lambda$.
		\begin{align*}
			\binom{n}{k} p_n^k (1-p_n)^{n-k} &= \frac{n!}{k!(n-k)!} \cdot \frac{\lambda_n^k}{n^k} \cdot \left(1 - \frac{\lambda_n}{n}\right)^{n-k} \\
			&= \frac{\lambda_n^k}{k!} \cdot \underbrace{\frac{n(n-1)\cdots(n-k+1)}{n^k}}_{\to 1} \cdot \underbrace{\left(1 - \frac{\lambda_n}{n}\right)^n}_{\to e^{-\lambda}} \cdot \underbrace{\left(1 - \frac{\lambda_n}{n}\right)^{-k}}_{\to 1}.
		\end{align*}
		
		Первый множитель $\to \lambda^k / k!$ (т.к. $\lambda_n \to \lambda$).
		
		Второй: $\dfrac{n(n-1)\cdots(n-k+1)}{n^k} = 1 \cdot \left(1 - \frac{1}{n}\right) \cdots \left(1 - \frac{k-1}{n}\right) \to 1$.
		
		Третий: $\left(1 - \dfrac{\lambda_n}{n}\right)^n \to e^{-\lambda}$ (стандартный предел).
		
		Итого: $\to \dfrac{\lambda^k}{k!} \cdot 1 \cdot e^{-\lambda} \cdot 1 = \dfrac{\lambda^k e^{-\lambda}}{k!}$.
	\end{proof}
	
	\begin{theorem}
		$\mathbb{E}[X] = \lambda$, $\mathbb{D}[X] = \lambda$ для $X \sim \text{Pois}(\lambda)$.
	\end{theorem}
	
	\begin{proof}
		\textit{Математическое ожидание:}
		\[\mathbb{E}[X] = \sum_{k=0}^{\infty} k \cdot \frac{\lambda^k e^{-\lambda}}{k!} = e^{-\lambda} \sum_{k=1}^{\infty} \frac{\lambda^k}{(k-1)!} = e^{-\lambda} \lambda \sum_{k=1}^{\infty} \frac{\lambda^{k-1}}{(k-1)!} = e^{-\lambda} \lambda e^{\lambda} = \lambda.\]
		
		\textit{Второй момент} (через $\mathbb{E}[X(X-1)]$):
		\[\mathbb{E}[X(X-1)] = \sum_{k=0}^{\infty} k(k-1) \cdot \frac{\lambda^k e^{-\lambda}}{k!} = e^{-\lambda} \sum_{k=2}^{\infty} \frac{\lambda^k}{(k-2)!} = e^{-\lambda} \lambda^2 e^{\lambda} = \lambda^2.\]
		
		Тогда $\mathbb{E}[X^2] = \mathbb{E}[X(X-1)] + \mathbb{E}[X] = \lambda^2 + \lambda$.
		
		\[\mathbb{D}[X] = \mathbb{E}[X^2] - (\mathbb{E}[X])^2 = \lambda^2 + \lambda - \lambda^2 = \lambda.\]
	\end{proof}
	
	\subsection{Неравенство Чебышёва и закон больших чисел}
	
	\begin{theorem}[Неравенство Маркова]
		Если $X \geqslant 0$ и $a > 0$, то $P(X \geqslant a) \leqslant \dfrac{\mathbb{E}[X]}{a}$.
	\end{theorem}
	
	\begin{proof}
		\[\mathbb{E}[X] = \sum_x x P(X = x) = \sum_{x < a} x P(X = x) + \sum_{x \geqslant a} x P(X = x) \geqslant \sum_{x \geqslant a} x P(X = x) \geqslant a \sum_{x \geqslant a} P(X = x) = a P(X \geqslant a).\]
	\end{proof}
	
	\begin{theorem}[Неравенство Чебышёва]
		Для любого $\varepsilon > 0$:
		\[P(|X - \mathbb{E}[X]| \geqslant \varepsilon) \leqslant \frac{\mathbb{D}[X]}{\varepsilon^2}.\]
	\end{theorem}
	
	\begin{proof}
		Применим неравенство Маркова к неотрицательной случайной величине $(X - \mathbb{E}[X])^2$ с $a = \varepsilon^2$:
		\[P(|X - \mathbb{E}[X]| \geqslant \varepsilon) = P\!\left((X - \mathbb{E}[X])^2 \geqslant \varepsilon^2\right) \leqslant \frac{\mathbb{E}[(X - \mathbb{E}[X])^2]}{\varepsilon^2} = \frac{\mathbb{D}[X]}{\varepsilon^2}.\]
	\end{proof}
	
	\begin{theorem}[Закон больших чисел (Чебышёв)]
		Пусть $X_1, X_2, \dots$~--- независимые одинаково распределённые случайные величины с $\mathbb{E}[X_i] = \mu$ и $\mathbb{D}[X_i] = \sigma^2 < \infty$. Тогда для любого $\varepsilon > 0$:
		\[P\!\left(\left|\frac{X_1 + \dots + X_n}{n} - \mu\right| \geqslant \varepsilon\right) \xrightarrow{n \to \infty} 0.\]
	\end{theorem}
	
	\begin{proof}
		Пусть $\overline{X}_n = \dfrac{X_1 + \dots + X_n}{n}$.
		
		По линейности: $\mathbb{E}[\overline{X}_n] = \dfrac{1}{n} \sum_{i=1}^n \mathbb{E}[X_i] = \mu$.
		
		По независимости: $\mathbb{D}[\overline{X}_n] = \mathbb{D}\!\left[\dfrac{1}{n}\sum_{i=1}^n X_i\right] = \dfrac{1}{n^2} \sum_{i=1}^n \mathbb{D}[X_i] = \dfrac{n\sigma^2}{n^2} = \dfrac{\sigma^2}{n}$.
		
		По неравенству Чебышёва:
		\[P(|\overline{X}_n - \mu| \geqslant \varepsilon) \leqslant \frac{\mathbb{D}[\overline{X}_n]}{\varepsilon^2} = \frac{\sigma^2}{n\varepsilon^2} \xrightarrow{n \to \infty} 0.\]
	\end{proof}
	
	\newpage
	
	\section{Ряды Фарея и близкие дроби. Цепные дроби, связь с алгоритмом Евклида. Формулы для подходящих дробей. Связь с близкими дробями. Порядок приближения иррациональных чисел подходящими дробями. Критерий периодичности цепных дробей.}
	
	\subsection{Ряды Фарея}
	
	\begin{definition}
		\textbf{Ряд Фарея} $F_n$~--- возрастающая последовательность всех несократимых дробей в $[0, 1]$ со знаменателями $\leqslant n$.
	\end{definition}
	
	Пример: $F_4 = \left\{\dfrac{0}{1}, \dfrac{1}{4}, \dfrac{1}{3}, \dfrac{1}{2}, \dfrac{2}{3}, \dfrac{3}{4}, \dfrac{1}{1}\right\}$.
	
	\begin{theorem}[Свойства ряда Фарея]
		\begin{enumerate}
			\item Если $\dfrac{a}{b}$ и $\dfrac{c}{d}$~--- соседние дроби в $F_n$, то $bc - ad = 1$.
			\item Если $\dfrac{a}{b}$, $\dfrac{c}{d}$, $\dfrac{e}{f}$~--- три последовательных дроби в $F_n$, то $\dfrac{c}{d} = \dfrac{a+e}{b+f}$ (медианта).
		\end{enumerate}
	\end{theorem}
	
	\begin{proof}[Доказательство пункта 1]
		Индукция по $n$. База $n = 1$: $F_1 = \{0/1, 1/1\}$, $bc - ad = 1 \cdot 1 - 0 \cdot 1 = 1$.
		
		Переход от $F_{n-1}$ к $F_n$: между соседними $\dfrac{a}{b}$ и $\dfrac{c}{d}$ в $F_{n-1}$ может вставиться новая дробь~--- медианта $\dfrac{a+c}{b+d}$~--- если $b + d = n$ (и медианта несократима).
		
		Проверяем для новых пар $\left(\dfrac{a}{b}, \dfrac{a+c}{b+d}\right)$ и $\left(\dfrac{a+c}{b+d}, \dfrac{c}{d}\right)$:
		\[b(a+c) - a(b+d) = bc - ad = 1,\]
		\[(a+c)d - c(b+d) = ad - bc = -(bc - ad) = -1... \]
		
		Стоп, должно быть $+1$. Пересчитаем: $c(b+d) - (a+c)d = cb + cd - ad - cd = bc - ad = 1$. Верно.
		
		Итого: свойство $bc - ad = 1$ сохраняется для всех соседних пар.
	\end{proof}
	
	\begin{corollary}
		Между соседними $\dfrac{a}{b}$ и $\dfrac{c}{d}$ в $F_n$ нет дробей со знаменателем $\leqslant n$: если $\dfrac{p}{q}$ строго между ними и $q \leqslant n$, то $|pc - qc| \geqslant 1$ и $|pa - qb| \geqslant 1$... (подробнее в связи с подходящими дробями).
	\end{corollary}
	
	\begin{definition}
		Дробь $\dfrac{p}{q}$ (в несократимом виде) называется \textbf{наилучшим приближением} (близкой дробью) к $\alpha$, если для любой дроби $\dfrac{r}{s}$ с $0 < s \leqslant q$ и $\dfrac{r}{s} \neq \dfrac{p}{q}$ выполнено:
		\[\left|\alpha - \frac{p}{q}\right| < \left|\alpha - \frac{r}{s}\right|.\]
	\end{definition}
	
	\subsection{Цепные дроби}
	
	\begin{definition}
		Конечная цепная дробь: $[a_0; a_1, a_2, \dots, a_n] = a_0 + \cfrac{1}{a_1 + \cfrac{1}{a_2 + \cfrac{1}{\ddots + \cfrac{1}{a_n}}}}$,
		
		где $a_0 \in \mathbb{Z}$, $a_i \in \mathbb{N}$ при $i \geqslant 1$ (называются \textbf{неполными частными}).
	\end{definition}
	
	\begin{definition}
		Бесконечная цепная дробь $[a_0; a_1, a_2, \dots] = \lim_{n \to \infty} [a_0; a_1, \dots, a_n]$.
	\end{definition}
	
	\begin{theorem}
		Каждое действительное число представляется в виде цепной дроби:
		\begin{itemize}
			\item рациональные~--- конечной цепной дробью (единственно, если $a_n > 1$);
			\item иррациональные~--- бесконечной цепной дробью (единственно).
		\end{itemize}
	\end{theorem}
	
	\subsection{Связь с алгоритмом Евклида}
	
	Для нахождения разложения $\dfrac{a}{b}$ в цепную дробь применяем алгоритм Евклида:
	\begin{align*}
		a &= b \cdot a_0 + r_1, \quad 0 \leqslant r_1 < b, \quad a_0 = \lfloor a/b \rfloor, \\
		b &= r_1 \cdot a_1 + r_2, \quad 0 \leqslant r_2 < r_1, \\
		r_1 &= r_2 \cdot a_2 + r_3, \\
		&\vdots \\
		r_{n-2} &= r_{n-1} \cdot a_{n-1} + r_n, \\
		r_{n-1} &= r_n \cdot a_n + 0.
	\end{align*}
	
	Тогда $\dfrac{a}{b} = [a_0; a_1, a_2, \dots, a_n]$.
	
	\begin{example}
		$\dfrac{17}{5}$: $17 = 5 \cdot 3 + 2$, $5 = 2 \cdot 2 + 1$, $2 = 1 \cdot 2 + 0$. Получаем $[3; 2, 2]$.
		
		Проверка: $3 + \cfrac{1}{2 + \cfrac{1}{2}} = 3 + \cfrac{1}{5/2} = 3 + \dfrac{2}{5} = \dfrac{17}{5}$. $\checkmark$
	\end{example}
	
	Для иррациональных чисел процесс не заканчивается. Например:
	\[\sqrt{2} = [1; 2, 2, 2, \dots] = [1; \overline{2}].\]
	
	\subsection{Подходящие дроби}
	
	\begin{definition}
		\textbf{Подходящие дроби (конвергенты):} $\dfrac{p_n}{q_n} = [a_0; a_1, \dots, a_n]$.
	\end{definition}
	
	\begin{theorem}[Рекуррентные формулы]
		\begin{align*}
			p_{-1} = 1, \quad p_0 &= a_0, \quad p_n = a_n p_{n-1} + p_{n-2} \text{ при } n \geqslant 1, \\
			q_{-1} = 0, \quad q_0 &= 1, \quad q_n = a_n q_{n-1} + q_{n-2} \text{ при } n \geqslant 1.
		\end{align*}
	\end{theorem}
	
	\begin{proof}
		Индукция по $n$.
		
		База $n = 0$: $p_0/q_0 = a_0/1$ и $a_0 \cdot 1 + 1 / (a_0 \cdot 0 + 1) = a_0$. $\checkmark$
		
		База $n = 1$: $[a_0; a_1] = a_0 + 1/a_1 = (a_0 a_1 + 1)/a_1$. Формула даёт $p_1 = a_1 a_0 + 1$, $q_1 = a_1 \cdot 1 + 0 = a_1$. $\checkmark$
		
		Переход: $[a_0; a_1, \dots, a_{n-1}, a_n] = [a_0; a_1, \dots, a_{n-2}, a_{n-1} + 1/a_n]$.
		
		Применяем формулу для $n - 1$ частных, но с последним элементом $a_{n-1}' = a_{n-1} + 1/a_n$:
		\begin{align*}
			\frac{p_n}{q_n} &= \frac{a_{n-1}' \cdot p_{n-2} + p_{n-3}}{a_{n-1}' \cdot q_{n-2} + q_{n-3}} = \frac{(a_{n-1} + 1/a_n) p_{n-2} + p_{n-3}}{(a_{n-1} + 1/a_n) q_{n-2} + q_{n-3}}.
		\end{align*}
		
		Умножаем числитель и знаменатель на $a_n$:
		\begin{align*}
			\frac{p_n}{q_n} &= \frac{a_n(a_{n-1} p_{n-2} + p_{n-3}) + p_{n-2}}{a_n(a_{n-1} q_{n-2} + q_{n-3}) + q_{n-2}} = \frac{a_n p_{n-1} + p_{n-2}}{a_n q_{n-1} + q_{n-2}}.
		\end{align*}
		
		Это и есть утверждение формулы.
	\end{proof}
	
	\begin{theorem}
		$p_n q_{n-1} - p_{n-1} q_n = (-1)^{n+1}$ для $n \geqslant 0$.
	\end{theorem}
	
	\begin{proof}
		Обозначим $\Delta_n = p_n q_{n-1} - p_{n-1} q_n$.
		
		\[\Delta_n = p_n q_{n-1} - p_{n-1} q_n = (a_n p_{n-1} + p_{n-2}) q_{n-1} - p_{n-1}(a_n q_{n-1} + q_{n-2}) = p_{n-2} q_{n-1} - p_{n-1} q_{n-2} = -\Delta_{n-1}.\]
		
		База: $\Delta_0 = p_0 q_{-1} - p_{-1} q_0 = a_0 \cdot 0 - 1 \cdot 1 = -1 = (-1)^1$.
		
		По рекуррентности: $\Delta_n = (-1)^{n+1}$.
	\end{proof}
	
	\begin{corollary}
		$\gcd(p_n, q_n) = 1$~--- подходящие дроби несократимы. (Из $\Delta_n = (-1)^{n+1}$ следует, что $\gcd(p_n, q_n)$ делит $1$.)
	\end{corollary}
	
	\begin{corollary}
		Разность последовательных подходящих дробей:
		\[\frac{p_n}{q_n} - \frac{p_{n-1}}{q_{n-1}} = \frac{p_n q_{n-1} - p_{n-1} q_n}{q_n q_{n-1}} = \frac{(-1)^{n+1}}{q_n q_{n-1}}.\]
		
		То есть подходящие дроби чётного порядка возрастают, нечётного~--- убывают, и все они «стягиваются» к пределу $\alpha = [a_0; a_1, a_2, \dots]$.
	\end{corollary}
	
	\begin{theorem}
		Для иррационального $\alpha = [a_0; a_1, a_2, \dots]$ и $n \geqslant 0$:
		\[\alpha = \frac{\alpha_{n+1} p_n + p_{n-1}}{\alpha_{n+1} q_n + q_{n-1}},\]
		где $\alpha_{n+1} = [a_{n+1}; a_{n+2}, \dots]$~--- $(n+1)$-е полное частное.
	\end{theorem}
	
	\begin{proof}
		По построению: $[a_0; a_1, \dots, a_n, \alpha_{n+1}]$, применяем рекуррентные формулы с «последним элементом» $\alpha_{n+1}$.
	\end{proof}
	
	\subsection{Порядок приближения}
	
	\begin{theorem}
		Для иррационального $\alpha$ и подходящей дроби $p_n/q_n$:
		\[\frac{1}{q_n(q_n + q_{n+1})} < \left|\alpha - \frac{p_n}{q_n}\right| < \frac{1}{q_n q_{n+1}}.\]
	\end{theorem}
	
	\begin{proof}
		Из формулы:
		\[\alpha - \frac{p_n}{q_n} = \frac{\alpha_{n+1} p_n + p_{n-1}}{\alpha_{n+1} q_n + q_{n-1}} - \frac{p_n}{q_n} = \frac{p_{n-1} q_n - p_n q_{n-1}}{q_n(\alpha_{n+1} q_n + q_{n-1})} = \frac{(-1)^n}{q_n(\alpha_{n+1} q_n + q_{n-1})}.\]
		
		Поскольку $\alpha_{n+1} = a_{n+1} + 1/\alpha_{n+2}$ и $a_{n+1} < \alpha_{n+1} < a_{n+1} + 1$:
		\[q_{n+1} = a_{n+1} q_n + q_{n-1} < \alpha_{n+1} q_n + q_{n-1} < (a_{n+1}+1) q_n + q_{n-1} = q_{n+1} + q_n.\]
		
		Значит:
		\[\frac{1}{q_n(q_{n+1} + q_n)} < \left|\alpha - \frac{p_n}{q_n}\right| < \frac{1}{q_n q_{n+1}}.\]
	\end{proof}
	
	\begin{corollary}
		Подходящие дроби хорошо приближают $\alpha$: $\left|\alpha - \dfrac{p_n}{q_n}\right| < \dfrac{1}{q_n q_{n+1}} < \dfrac{1}{q_n^2}$.
	\end{corollary}
	
	\begin{theorem}[Теорема Лежандра]
		Если $\left|\alpha - \dfrac{p}{q}\right| < \dfrac{1}{2q^2}$, то $\dfrac{p}{q}$~--- подходящая дробь к $\alpha$.
	\end{theorem}
	
	\begin{proof}
		Предположим $p/q$ не является подходящей дробью. Тогда существует $n$ такое, что $q_n \leqslant q < q_{n+1}$.
		
		Подходящая дробь $p_n/q_n$ является лучшим приближением с знаменателем $\leqslant q$ (по свойству подходящих дробей), то есть:
		\[\left|\alpha - \frac{p_n}{q_n}\right| \leqslant \left|\alpha - \frac{p}{q}\right| < \frac{1}{2q^2} \leqslant \frac{1}{2q_n q}.\]
		
		С другой стороны, $p/q \neq p_n/q_n$~--- дроби с разными знаменателями (или разные дроби), поэтому:
		\[\left|\frac{p}{q} - \frac{p_n}{q_n}\right| = \frac{|p q_n - p_n q|}{q q_n} \geqslant \frac{1}{q q_n}.\]
		
		По неравенству треугольника:
		\[\frac{1}{q q_n} \leqslant \left|\frac{p}{q} - \frac{p_n}{q_n}\right| \leqslant \left|\frac{p}{q} - \alpha\right| + \left|\alpha - \frac{p_n}{q_n}\right| < \frac{1}{2q^2} + \frac{1}{2q_n q} \leqslant \frac{1}{q_n q}.\]
		
		Получили $\dfrac{1}{q q_n} < \dfrac{1}{q_n q}$~--- противоречие. Значит, $p/q$ является подходящей дробью.
	\end{proof}
	
	\subsection{Связь с близкими дробями}
	
	\begin{theorem}
		Всякая подходящая дробь является наилучшим приближением первого рода (то есть близкой дробью).
	\end{theorem}
	
	\begin{proof}
		Пусть $p/q$~--- дробь с $q \leqslant q_n$ и $p/q \neq p_n/q_n$. Покажем $|\alpha - p_n/q_n| < |\alpha - p/q|$.
		
		Пусть $p_n/q_n$~--- $n$-я подходящая дробь, $p_{n-1}/q_{n-1}$~--- $(n-1)$-я. Рассмотрим систему:
		\[p = c p_n + d p_{n-1}, \quad q = c q_n + d q_{n-1}.\]
		
		Из $p_n q_{n-1} - p_{n-1} q_n = (-1)^{n+1}$: $c = (-1)^{n+1}(p q_{n-1} - p_{n-1} q)$, $d = (-1)^n(p_n q - p q_n)$.
		
		Если $d = 0$, то $p/q = p_n/q_n$~--- противоречие. Значит $d \neq 0$.
		
		$\alpha - p/q = c(\alpha - p_n/q_n) + d(\alpha - p_{n-1}/q_{n-1})$.
		
		Поскольку $\alpha - p_n/q_n$ и $\alpha - p_{n-1}/q_{n-1}$ имеют разные знаки (подходящие дроби с двух сторон), и $|c|, |d| \geqslant 0$ с $d \neq 0$:
		\[|\alpha - p/q| \geqslant |d| \cdot |\alpha - p_{n-1}/q_{n-1}| > |\alpha - p_n/q_n|.\]
	\end{proof}
	
	\begin{theorem}
		Всякое наилучшее приближение является подходящей или промежуточной дробью $\dfrac{cp_n + p_{n-1}}{cq_n + q_{n-1}}$ для $1 \leqslant c \leqslant a_{n+1}$.
	\end{theorem}
	
	\subsection{Критерий периодичности}
	
	\begin{definition}
		Цепная дробь $[a_0; a_1, a_2, \dots]$ \textbf{периодична}, если существуют $k \geqslant 0$ и $T > 0$ такие, что $a_{n+T} = a_n$ для всех $n \geqslant k$. Обозначение: $[a_0; a_1, \dots, a_{k-1}, \overline{a_k, \dots, a_{k+T-1}}]$.
	\end{definition}
	
	\begin{theorem}[Лагранж]
		Цепная дробь числа $\alpha$ периодична тогда и только тогда, когда $\alpha$~--- квадратичная иррациональность, то есть корень уравнения $ax^2 + bx + c = 0$ с $a, b, c \in \mathbb{Z}$, $a \neq 0$, $b^2 - 4ac > 0$, $b^2 - 4ac$ не является полным квадратом.
	\end{theorem}
	
	\begin{proof}
		\textbf{($\Leftarrow$) Если $\alpha$ квадратично иррационально, то цепная дробь периодична.}
		
		Запишем $\alpha = \dfrac{P + \sqrt{D}}{Q}$, где $P, Q, D \in \mathbb{Z}$, $D > 0$, $D$ не квадрат, $Q \mid D - P^2$ (можно привести к такому виду).
		
		Полные частные определяются как $\alpha_0 = \alpha$, $\alpha_{n+1} = \dfrac{1}{\alpha_n - a_n}$, где $a_n = \lfloor \alpha_n \rfloor$.
		
		Утверждение: $\alpha_n = \dfrac{P_n + \sqrt{D}}{Q_n}$ для некоторых $P_n, Q_n \in \mathbb{Z}$ с $Q_n \mid D - P_n^2$.
		
		Это верно для $n = 0$. Если $\alpha_n = \dfrac{P_n + \sqrt{D}}{Q_n}$, то:
		\[\alpha_{n+1} = \frac{1}{\alpha_n - a_n} = \frac{Q_n}{P_n - a_n Q_n + \sqrt{D}} = \frac{Q_n(-(P_n - a_n Q_n) + \sqrt{D})}{D - (P_n - a_n Q_n)^2}.\]
		
		Полагая $P_{n+1} = a_n Q_n - P_n$ и $Q_{n+1} = \dfrac{D - P_{n+1}^2}{Q_n}$ (целое, т.к. $Q_n \mid D - P_n^2$ и ... индукция), получаем $\alpha_{n+1} = \dfrac{P_{n+1} + \sqrt{D}}{Q_{n+1}}$.
		
		Можно показать, что $|P_n| \leqslant \sqrt{D}$ и $|Q_n| \leqslant 2\sqrt{D}$ при $n \geqslant 1$ (оценки из условия $0 < \alpha_n - a_n < 1$). Поскольку $P_n, Q_n$ целые и ограничены, множество значений $(P_n, Q_n)$ конечно. Значит последовательность $(\alpha_n)$ периодична~--- цепная дробь периодична.
		
		\textbf{($\Rightarrow$) Если цепная дробь периодична, то $\alpha$ квадратично иррационально.}
		
		Пусть $\alpha = [a_0; \dots, a_{k-1}, \overline{a_k, \dots, a_{k+T-1}}]$. Обозначим $\beta = [{\overline{a_k, \dots, a_{k+T-1}}}]$~--- чисто периодическая часть.
		
		Из периодичности: $\beta = [a_k; a_{k+1}, \dots, a_{k+T-1}, \beta]$ (подставляем $\beta$ в конец периода).
		
		Применяем рекуррентную формулу: $\beta = \dfrac{\beta \cdot p_{T-1} + p_{T-2}}{\beta \cdot q_{T-1} + q_{T-2}}$, где $p_i/q_i$~--- подходящие дроби для $[a_k; \dots, a_{k+T-1}]$.
		
		Это даёт квадратное уравнение $q_{T-1} \beta^2 + (q_{T-2} - p_{T-1})\beta - p_{T-2} = 0$ с целыми коэффициентами, откуда $\beta$~--- квадратично иррационально.
		
		Тогда $\alpha = [a_0; \dots, a_{k-1}, \beta] = \dfrac{\beta p_{k-1} + p_{k-2}}{\beta q_{k-1} + q_{k-2}}$~--- дробно-линейная функция от $\beta$, и тоже квадратично иррационально.
	\end{proof}
	
	\begin{example}
		$\sqrt{2} = [1; \overline{2}]$. Полное частное $\beta = [\overline{2}]$ удовлетворяет $\beta = 2 + 1/\beta$, то есть $\beta^2 - 2\beta - 1 = 0$, $\beta = 1 + \sqrt{2}$. Тогда $\sqrt{2} = [1; \beta] = -1 + \beta = \sqrt{2}$. $\checkmark$
	\end{example}
	
	\newpage
	
	\section{Понятие группы, примеры групп. Подгруппы, смежные классы, теорема Лагранжа. Нормальные подгруппы, эквивалентность определений. Факторгруппа. Гомоморфизм, ядро, образ. Теоремы о гомоморфизме. Перестановки, чётность перестановок, критерий сопряжённости перестановок. Группа автоморфизмов графа, теорема Фрухта. Классификация конечных абелевых групп (без доказательства).}
	
	\subsection{Понятие группы}
	
	\begin{definition}
		\textbf{Группа}~--- множество $G$ с бинарной операцией $\cdot: G \times G \to G$, удовлетворяющей:
		\begin{enumerate}
			\item (Ассоциативность) $(ab)c = a(bc)$ для всех $a, b, c \in G$
			\item (Нейтральный элемент) $\exists e \in G: ae = ea = a$ для всех $a \in G$
			\item (Обратный элемент) $\forall a \in G \ \exists a^{-1} \in G: aa^{-1} = a^{-1}a = e$
		\end{enumerate}
		Группа \textbf{абелева}, если дополнительно $ab = ba$ для всех $a, b \in G$.
	\end{definition}
	
	\begin{note}
		Нейтральный и обратный элементы единственны.
		
		\textit{Единственность нейтрального:} если $e$ и $e'$~--- оба нейтральных, то $e = e \cdot e' = e'$.
		
		\textit{Единственность обратного:} если $ba = ab = e$ и $ca = ac = e$, то $b = b(ac) = (ba)c = c$.
	\end{note}
	
	\textbf{Примеры групп:}
	\begin{itemize}
		\item $(\mathbb{Z}, +)$, $(\mathbb{Q}, +)$, $(\mathbb{R}, +)$~--- абелевы группы.
		\item $(\mathbb{Q}^*, \cdot)$, $(\mathbb{R}^*, \cdot)$, $(\mathbb{C}^*, \cdot)$~--- абелевы группы.
		\item $(\mathbb{Z}_n, +)$~--- вычеты по модулю $n$, $n$ элементов.
		\item $(\mathbb{Z}_n^*, \cdot)$~--- вычеты, взаимно простые с $n$, $\phi(n)$ элементов.
		\item $S_n$~--- группа перестановок $\{1, \dots, n\}$, $|S_n| = n!$ (неабелева при $n \geqslant 3$).
		\item $GL_n(\mathbb{R}) = \{A \in M_n(\mathbb{R}) : \det A \neq 0\}$~--- обратимые матрицы.
		\item $D_n$~--- группа симметрий правильного $n$-угольника, $|D_n| = 2n$ (вращения и отражения).
		\item Тривиальная группа $\{e\}$.
	\end{itemize}
	
	\subsection{Подгруппы и смежные классы}
	
	\begin{definition}
		$H \subseteq G$~--- \textbf{подгруппа} ($H \leqslant G$), если $H$ замкнута относительно операции $G$ и сама является группой.
	\end{definition}
	
	\begin{theorem}[Критерий подгруппы]
		$H \leqslant G \Leftrightarrow$ (1) $H \neq \varnothing$ и (2) $\forall a, b \in H: ab^{-1} \in H$.
	\end{theorem}
	
	\begin{proof}
		$(\Rightarrow)$ Очевидно из аксиом группы.
		
		$(\Leftarrow)$ Из $H \neq \varnothing$ возьмём $a \in H$. Тогда $aa^{-1} = e \in H$. Из $ea^{-1} = a^{-1} \in H$. Для $a, b \in H$: $a \cdot (b^{-1})^{-1} = ab \in H$.
		Ассоциативность наследуется от $G$.
	\end{proof}
	
	\begin{definition}
		\textbf{Левый смежный класс} элемента $a$ по подгруппе $H$: $aH = \{ah : h \in H\}$.
		
		\textbf{Правый смежный класс}: $Ha = \{ha : h \in H\}$.
	\end{definition}
	
	\begin{lemma}
		Левые смежные классы образуют разбиение $G$.
	\end{lemma}
	
	\begin{proof}
		$a \in aH$ (т.к. $e \in H$), поэтому классы покрывают $G$. Если $aH \cap bH \neq \varnothing$, то $\exists h_1, h_2 \in H: ah_1 = bh_2$, откуда $a = bh_2 h_1^{-1} \in bH$. Тогда $aH = bH$ (для любого $ah \in aH$: $ah = b(h_2 h_1^{-1} h) \in bH$, и обратно).
	\end{proof}
	
	\begin{theorem}[Лагранж]
		Если $G$ конечна и $H \leqslant G$, то $|H|$ делит $|G|$. Точнее, $|G| = [G : H] \cdot |H|$.
	\end{theorem}
	
	\begin{proof}
		Левые смежные классы образуют разбиение $G$. Все смежные классы равномощны: биекция $H \to aH$ задаётся $h \mapsto ah$ (обратная: $ah \mapsto h$). Если число классов равно $k = [G : H]$, то $|G| = k \cdot |H|$.
	\end{proof}
	
	\begin{definition}
		\textbf{Индекс подгруппы} $[G : H]$~--- число левых смежных классов $G$ по $H$.
	\end{definition}
	
	\begin{corollary}
		Порядок любого элемента $g \in G$ делит $|G|$ (порядок элемента~--- это $|\langle g \rangle|$, где $\langle g \rangle = \{e, g, g^2, \dots\}$~--- подгруппа).
	\end{corollary}
	
	\subsection{Нормальные подгруппы}
	
	\begin{definition}
		$H \trianglelefteq G$ (\textbf{нормальная подгруппа}), если выполнено одно из эквивалентных условий:
		\begin{enumerate}
			\item $\forall g \in G: gH = Hg$ (левые и правые смежные классы совпадают)
			\item $\forall g \in G: gHg^{-1} = H$ (инвариантность при сопряжении)
			\item $\forall g \in G, h \in H: ghg^{-1} \in H$ (замкнутость относительно сопряжений)
		\end{enumerate}
	\end{definition}
	
	\begin{proof}[Эквивалентность определений]
		$(1) \Rightarrow (2)$: $gH = Hg \Rightarrow gHg^{-1} = Hgg^{-1} = He = H$.
		
		$(2) \Rightarrow (3)$: если $gHg^{-1} = H$, то для $h \in H$: $ghg^{-1} \in gHg^{-1} = H$.
		
		$(3) \Rightarrow (1)$: Покажем $gH \subseteq Hg$: для $h \in H$ имеем $gh = (ghg^{-1})g \in Hg$ (т.к. $ghg^{-1} \in H$ по (3)). Аналогично $Hg \subseteq gH$ (применяем (3) к $g^{-1}$: $g^{-1}h g \in H$, значит $hg = g(g^{-1}hg) \in gH$).
	\end{proof}
	
	\begin{example}
		\begin{itemize}
			\item $\{e\}$ и $G$ всегда нормальны.
			\item Подгруппа индекса 2 всегда нормальна: два смежных класса, один из них $H$, другой $G \setminus H$~--- одинаковы слева и справа.
			\item $A_n \trianglelefteq S_n$ (знакочередующая группа, индекс 2).
			\item Центр $Z(G) = \{g \in G : gx = xg \ \forall x \in G\} \trianglelefteq G$.
		\end{itemize}
	\end{example}
	
	\subsection{Факторгруппа}
	
	\begin{theorem}
		Если $H \trianglelefteq G$, то $G/H = \{gH : g \in G\}$ с операцией $(aH)(bH) = (ab)H$~--- группа (\textbf{факторгруппа}).
	\end{theorem}
	
	\begin{proof}
		\textit{Корректность операции:} если $aH = a'H$ и $bH = b'H$, то $a' = ah_1$, $b' = bh_2$ для некоторых $h_1, h_2 \in H$. Тогда:
		\[a'b' = ah_1 \cdot bh_2 = a \cdot (h_1 b) \cdot h_2.\]
		Так как $H \trianglelefteq G$: $h_1 b = b h_3$ для некоторого $h_3 \in H$ (из $bH = Hb$). Значит $a'b' = ab \cdot h_3 h_2 \in (ab)H$, откуда $a'b'H = abH$.
		
		\textit{Ассоциативность:} $((aH)(bH))(cH) = (abH)(cH) = (ab)cH = a(bc)H = (aH)((bc)H) = (aH)((bH)(cH))$.
		
		\textit{Нейтральный:} $eH = H$; $(aH)(eH) = aeH = aH$.
		
		\textit{Обратный:} $(aH)(a^{-1}H) = aa^{-1}H = eH = H$.
	\end{proof}
	
	\begin{note}
		Естественная проекция $\pi: G \to G/H$, $\pi(g) = gH$~--- гомоморфизм: $\pi(ab) = abH = (aH)(bH) = \pi(a)\pi(b)$.
	\end{note}
	
	\subsection{Гомоморфизм}
	
	\begin{definition}
		Отображение $\varphi: G \to H$~--- \textbf{гомоморфизм групп}, если $\varphi(ab) = \varphi(a)\varphi(b)$ для всех $a, b \in G$.
	\end{definition}
	
	\begin{note}
		Из определения следует: $\varphi(e_G) = e_H$ и $\varphi(a^{-1}) = \varphi(a)^{-1}$.
		
		\textit{Доказательство:} $\varphi(e_G) \cdot \varphi(e_G) = \varphi(e_G \cdot e_G) = \varphi(e_G)$, умножаем на $\varphi(e_G)^{-1}$: $\varphi(e_G) = e_H$. Далее: $\varphi(a) \cdot \varphi(a^{-1}) = \varphi(a \cdot a^{-1}) = \varphi(e_G) = e_H$.
	\end{note}
	
	\begin{definition}
		\textbf{Ядро:} $\ker \varphi = \{g \in G : \varphi(g) = e_H\}$.
		
		\textbf{Образ:} $\operatorname{Im} \varphi = \varphi(G) = \{\varphi(g) : g \in G\} \subseteq H$.
	\end{definition}
	
	\begin{theorem}
		$\ker \varphi \trianglelefteq G$, $\operatorname{Im} \varphi \leqslant H$.
		
		Кроме того, $\varphi$ инъективно $\Leftrightarrow \ker \varphi = \{e_G\}$.
	\end{theorem}
	
	\begin{proof}
		\textit{$\ker \varphi \leqslant G$:} $e_G \in \ker \varphi$, значит непусто. Если $a, b \in \ker \varphi$: $\varphi(ab^{-1}) = \varphi(a)\varphi(b)^{-1} = e_H \cdot e_H = e_H$.
		
		\textit{$\ker \varphi \trianglelefteq G$:} для $k \in \ker \varphi$, $g \in G$: $\varphi(gkg^{-1}) = \varphi(g)\varphi(k)\varphi(g)^{-1} = \varphi(g) \cdot e_H \cdot \varphi(g)^{-1} = e_H$.
		
		\textit{$\operatorname{Im} \varphi \leqslant H$:} $\varphi(e_G) = e_H \in \operatorname{Im} \varphi$. Если $\varphi(a), \varphi(b) \in \operatorname{Im} \varphi$: $\varphi(a)\varphi(b)^{-1} = \varphi(a)\varphi(b^{-1}) = \varphi(ab^{-1}) \in \operatorname{Im} \varphi$.
		
		\textit{Инъективность:} если $\ker \varphi = \{e\}$ и $\varphi(a) = \varphi(b)$, то $\varphi(ab^{-1}) = \varphi(a)\varphi(b)^{-1} = e_H$, значит $ab^{-1} = e$, то есть $a = b$. Обратно: если $\varphi$ инъективно и $\varphi(a) = e_H = \varphi(e)$, то $a = e$.
	\end{proof}
	
	\begin{theorem}[Первая теорема о гомоморфизме]
		$G / \ker \varphi \cong \operatorname{Im} \varphi$.
	\end{theorem}
	
	\begin{proof}
		Пусть $K = \ker \varphi$. Определим $\psi: G/K \to \operatorname{Im}\varphi$: $\psi(gK) = \varphi(g)$.
		
		\textit{Корректность:} если $gK = g'K$, то $g^{-1}g' \in K$, значит $\varphi(g^{-1}g') = e_H$, откуда $\varphi(g)^{-1}\varphi(g') = e_H$, то есть $\varphi(g) = \varphi(g')$.
		
		\textit{Гомоморфизм:} $\psi(aK \cdot bK) = \psi(abK) = \varphi(ab) = \varphi(a)\varphi(b) = \psi(aK)\psi(bK)$.
		
		\textit{Инъективность:} если $\psi(gK) = e_H$, то $\varphi(g) = e_H$, значит $g \in K$, то есть $gK = K = e_{G/K}$.
		
		\textit{Сюръективность на $\operatorname{Im}\varphi$:} для любого $\varphi(g) \in \operatorname{Im}\varphi$ имеем $\psi(gK) = \varphi(g)$.
	\end{proof}
	
	\begin{theorem}[Вторая теорема о гомоморфизме]
		Пусть $H \leqslant G$, $N \trianglelefteq G$. Тогда $HN \leqslant G$, $H \cap N \trianglelefteq H$ и $H/(H \cap N) \cong HN/N$.
	\end{theorem}
	
	\begin{proof}
		Рассмотрим ограничение проекции $\pi: G \to G/N$ на $H$: $\varphi = \pi|_H: H \to G/N$, $\varphi(h) = hN$.
		
		$\ker \varphi = \{h \in H : hN = N\} = \{h \in H : h \in N\} = H \cap N$.
		
		$\operatorname{Im} \varphi = \{hN : h \in H\} = HN/N$.
		
		По первой теореме: $H/(H \cap N) \cong HN/N$.
	\end{proof}
	
	\subsection{Перестановки}
	
	\begin{definition}
		\textbf{Перестановка}~--- биекция $\sigma: \{1, \dots, n\} \to \{1, \dots, n\}$.
		$S_n$~--- группа всех перестановок (относительно композиции), $|S_n| = n!$.
	\end{definition}
	
	\begin{definition}
		\textbf{Цикл} $(i_1 \, i_2 \, \dots \, i_k)$~--- перестановка: $i_1 \mapsto i_2 \mapsto \dots \mapsto i_k \mapsto i_1$, остальные элементы неподвижны.
		
		\textbf{Длина цикла}~--- число $k$ переставляемых элементов.
		
		\textbf{Транспозиция}~--- цикл длины 2.
		
		Два цикла \textbf{независимы (дизъюнктны)}, если они не имеют общих элементов; такие циклы коммутируют.
	\end{definition}
	
	\begin{theorem}
		Любая перестановка раскладывается в произведение попарно независимых циклов единственным образом (с точностью до порядка сомножителей и выбора начального элемента цикла).
	\end{theorem}
	
	\begin{proof}
		\textit{Существование:} для каждого элемента $i$ рассмотрим его орбиту $\{i, \sigma(i), \sigma^2(i), \dots\}$~--- конечное множество (так как $\{1, \dots, n\}$ конечно). Перестановка $\sigma$ действует на своих орбитах циклически. Это даёт разложение в произведение циклов.
		
		\textit{Единственность:} орбиты однозначно определяются перестановкой.
	\end{proof}
	
	\begin{definition}
		\textbf{Цикловой тип} перестановки~--- набор длин циклов в её разложении (обычно записывается в порядке убывания).
	\end{definition}
	
	\begin{theorem}
		Любая перестановка раскладывается в произведение транспозиций (не единственным образом, но чётность числа транспозиций сохраняется).
	\end{theorem}
	
	\begin{proof}
		$(i_1 \, i_2 \, \dots \, i_k) = (i_1 \, i_k)(i_1 \, i_{k-1}) \cdots (i_1 \, i_2)$.
		
		Проверка: $(i_1 \, i_2)(i_1 \, i_3) = (i_1 \, i_3 \, i_2)$... на самом деле, проверим общую формулу:
		$(i_1 \, i_k)\cdots(i_1 \, i_2)$: элемент $i_1 \mapsto i_2$ (последней транспозицией), $i_j \mapsto i_1 \mapsto i_{j+1}$ (транспозицией $(i_1 \, i_j)$, потом $(i_1 \, i_{j+1})$), $i_k \mapsto i_1$ (первой). $\checkmark$
	\end{proof}
	
	\begin{definition}
		\textbf{Знак (чётность) перестановки:}
		\[\operatorname{sgn}(\sigma) = (-1)^{N(\sigma)},\]
		где $N(\sigma)$ — число инверсий: $N(\sigma) = |\{(i,j) : i < j, \sigma(i) > \sigma(j)\}|$.
		
		Эквивалентно: $\operatorname{sgn}(\sigma) = \prod_{1 \leqslant i < j \leqslant n} \dfrac{\sigma(j) - \sigma(i)}{j - i}$.
	\end{definition}
	
	\begin{theorem}
		$\operatorname{sgn}$~--- гомоморфизм $S_n \to \{+1, -1\}$. В частности:
		\begin{enumerate}
			\item $\operatorname{sgn}(\sigma \tau) = \operatorname{sgn}(\sigma) \cdot \operatorname{sgn}(\tau)$
			\item Любая транспозиция имеет знак $-1$
			\item Цикл длины $k$ имеет знак $(-1)^{k-1}$
			\item Чётность числа транспозиций в любом разложении перестановки определена однозначно
		\end{enumerate}
	\end{theorem}
	
	\begin{proof}
		Из формулы $\operatorname{sgn}(\sigma) = \prod_{i < j} \dfrac{\sigma(j) - \sigma(i)}{j - i}$ следует $\operatorname{sgn}(\sigma\tau) = \operatorname{sgn}(\sigma)\operatorname{sgn}(\tau)$ (подстановка).
		
		Транспозиция $(ab)$ с $a < b$: меняет знак у пар $(a, b)$, $(a, c)$ для $a < c < b$ (меняет дважды~--- знак не меняется) и у пар $(c, b)$ для $a < c < b$ (тоже дважды), и у пары $(a, b)$ (один раз). Итого: $\operatorname{sgn}((ab)) = -1$.
		
		Цикл $(i_1 \, \dots \, i_k) = (i_1 \, i_k) \cdots (i_1 \, i_2)$~--- произведение $k - 1$ транспозиций, знак $(-1)^{k-1}$.
	\end{proof}
	
	\begin{definition}
		$A_n = \{\sigma \in S_n : \operatorname{sgn}(\sigma) = 1\}$~--- \textbf{знакочередующая группа} (альтернирующая).
	\end{definition}
	
	\begin{theorem}
		$A_n \trianglelefteq S_n$, $|A_n| = n!/2$, $[S_n : A_n] = 2$.
	\end{theorem}
	
	\begin{proof}
		$A_n = \ker \operatorname{sgn}$~--- ядро гомоморфизма, значит нормальная подгруппа. Гомоморфизм $\operatorname{sgn}: S_n \to \{+1, -1\}$ сюръективен (транспозиция имеет знак $-1$), поэтому по первой теореме $S_n / A_n \cong \{+1, -1\}$, откуда $[S_n : A_n] = 2$.
	\end{proof}
	
	\begin{theorem}[Критерий сопряжённости в $S_n$]
		$\sigma$ и $\tau$ сопряжены в $S_n$ (то есть $\exists \rho \in S_n: \tau = \rho \sigma \rho^{-1}$) $\Leftrightarrow$ они имеют одинаковый цикловой тип.
	\end{theorem}
	
	\begin{proof}
		$(\Rightarrow)$ Если $\sigma = (i_1 \, i_2 \, \dots)(j_1 \, j_2 \, \dots)\cdots$, то:
		\[\rho \sigma \rho^{-1} = (\rho(i_1) \, \rho(i_2) \, \dots)(\rho(j_1) \, \rho(j_2) \, \dots)\cdots\]
		
		Проверим: $(\rho \sigma \rho^{-1})(\rho(i_k)) = \rho(\sigma(i_k)) = \rho(i_{k+1})$. $\checkmark$
		
		Цикловой тип (набор длин) не меняется.
		
		$(\Leftarrow)$ Если $\sigma = (i_1 \, \dots \, i_{k_1})(i_{k_1+1} \, \dots \, i_{k_2})\cdots$ и $\tau = (j_1 \, \dots \, j_{k_1})(j_{k_1+1} \, \dots \, j_{k_2})\cdots$ имеют одинаковый цикловой тип, определим $\rho$ как перестановку, переводящую $i_m \mapsto j_m$ (по порядку внутри соответствующих циклов). Тогда $\rho \sigma \rho^{-1} = \tau$.
	\end{proof}
	
	\begin{corollary}
		Число классов сопряжённости $S_n$~--- это число разбиений числа $n$ (число способов представить $n$ как сумму натуральных чисел, с точностью до порядка).
	\end{corollary}
	
	\subsection{Автоморфизмы графа}
	
	\begin{definition}
		\textbf{Автоморфизм графа} $G = (V, E)$~--- биекция $\varphi: V \to V$, сохраняющая смежность: $\{u, v\} \in E \Leftrightarrow \{\varphi(u), \varphi(v)\} \in E$.
		
		$\operatorname{Aut}(G)$~--- группа автоморфизмов (относительно композиции).
	\end{definition}
	
	\begin{example}
		\begin{itemize}
			\item Полный граф $K_n$: $\operatorname{Aut}(K_n) = S_n$.
			\item Путь $P_n$: $\operatorname{Aut}(P_n) = \mathbb{Z}_2$ (единственный нетривиальный автоморфизм~--- «отражение»).
			\item Цикл $C_n$: $\operatorname{Aut}(C_n) = D_n$ (группа симметрий правильного $n$-угольника).
			\item Асимметричный (жёсткий) граф: $\operatorname{Aut}(G) = \{e\}$.
		\end{itemize}
	\end{example}
	
	\begin{theorem}[Фрухт, 1939]
		Для любой конечной группы $G$ существует конечный граф $\Gamma$ такой, что $\operatorname{Aut}(\Gamma) \cong G$.
	\end{theorem}
	
	\begin{proof}[Идея доказательства]
		По теореме Кэли, $G$ вкладывается в $S_{|G|}$ (как группа левых сдвигов). Граф Кэли $\Gamma$ группы $G$ по системе образующих $S$ имеет вершины $G$ и рёбра $\{g, gs\}$ для $g \in G$, $s \in S$.
		
		Граф Кэли имеет $G$ как подгруппу своего группы автоморфизмов. Чтобы «убить» лишние автоморфизмы, заменяем каждое ребро (помеченное образующей $s$) на «гаджет»~--- небольшой граф, однозначно идентифицирующий направление и тип ребра.
		
		Конкретно: каждому ребру $\{g, gs\}$ с образующей $s_i$ сопоставляется путь длины $n_i$ (где $n_i$ различны для разных образующих и их обратных). Это делает все рёбра «различимыми», и автоморфизмы полученного графа соответствуют ровно элементам $G$.
	\end{proof}
	
	\subsection{Классификация конечных абелевых групп}
	
	\begin{theorem}[Без доказательства]
		Любая конечная абелева группа изоморфна прямому произведению циклических групп простых степенных порядков:
		\[G \cong \mathbb{Z}_{p_1^{a_1}} \times \mathbb{Z}_{p_2^{a_2}} \times \cdots \times \mathbb{Z}_{p_k^{a_k}},\]
		где $p_i$~--- простые (не обязательно различные), $a_i \geqslant 1$. Разложение единственно с точностью до порядка сомножителей.
	\end{theorem}
	
	\begin{example}
		Все абелевы группы порядка 12: $12 = 4 \cdot 3 = 2^2 \cdot 3$.
		
		Возможные разложения:
		\[\mathbb{Z}_4 \times \mathbb{Z}_3 \cong \mathbb{Z}_{12}, \quad \mathbb{Z}_2 \times \mathbb{Z}_2 \times \mathbb{Z}_3.\]
		
		Групп ровно 2 (с точностью до изоморфизма).
	\end{example}
	
	\begin{note}
		Эквивалентная формулировка~--- через инвариантные делители: $G \cong \mathbb{Z}_{d_1} \times \cdots \times \mathbb{Z}_{d_r}$, где $d_1 \mid d_2 \mid \cdots \mid d_r$.
	\end{note}
	
	\newpage
	
	\section{Метрические пространства. Примеры. Непрерывные отображения метрических пространств, равносильные определения. Связность, теорема о промежуточном значении, задачи о блинах. Компактность, критерий компактности в $\mathbb{R}^n$. Теорема Вейерштрасса о достижении максимума.}
	
	\subsection{Метрические пространства}
	
	\begin{definition}
		\textbf{Метрическое пространство}~--- пара $(X, d)$, где $X$~--- множество, $d: X \times X \to [0, +\infty)$~--- метрика, удовлетворяющая:
		\begin{enumerate}
			\item (Невырожденность) $d(x, y) = 0 \Leftrightarrow x = y$
			\item (Симметричность) $d(x, y) = d(y, x)$
			\item (Неравенство треугольника) $d(x, z) \leqslant d(x, y) + d(y, z)$
		\end{enumerate}
	\end{definition}
	
	\subsection{Примеры}
	
	\textbf{1. $\mathbb{R}^n$ с евклидовой метрикой:} $d_2(x, y) = \sqrt{\sum_{i=1}^{n} (x_i - y_i)^2}$.
	
	\begin{proof}[Проверка неравенства треугольника]
		Неравенство треугольника~--- это неравенство Минковского:
		\[\sqrt{\sum_i (x_i - z_i)^2} \leqslant \sqrt{\sum_i (x_i - y_i)^2} + \sqrt{\sum_i (y_i - z_i)^2}.\]
		Это следует из неравенства Коши--Шварца: $\sum_i a_i b_i \leqslant \sqrt{\sum_i a_i^2} \cdot \sqrt{\sum_i b_i^2}$.
	\end{proof}
	
	\textbf{2. $\mathbb{R}^n$ с манхэттенской метрикой:} $d_1(x, y) = \sum_{i=1}^{n} |x_i - y_i|$.
	
	\begin{proof}
		Неравенство треугольника: $|x_i - z_i| \leqslant |x_i - y_i| + |y_i - z_i|$, суммируем по $i$.
	\end{proof}
	
	\textbf{3. $\mathbb{R}^n$ с метрикой Чебышёва:} $d_\infty(x, y) = \max_{1 \leqslant i \leqslant n} |x_i - y_i|$.
	
	\begin{proof}
		$|x_i - z_i| \leqslant |x_i - y_i| + |y_i - z_i| \leqslant d_\infty(x, y) + d_\infty(y, z)$ для каждого $i$. Берём максимум по $i$.
	\end{proof}
	
	\begin{note}
		Все три метрики на $\mathbb{R}^n$ эквивалентны: $d_\infty \leqslant d_2 \leqslant d_1 \leqslant n \cdot d_\infty$.
	\end{note}
	
	\textbf{4. Дискретная метрика:} $d(x, y) = \begin{cases} 0, & x = y \\ 1, & x \neq y \end{cases}$
	
	\begin{proof}
		Неравенство треугольника: если $x \neq z$, то $d(x,z) = 1$. Из $x \neq z$: либо $x \neq y$, либо $y \neq z$ (или оба), значит $d(x,y) + d(y,z) \geqslant 1$.
	\end{proof}
	
	\textbf{5. Пространство ограниченных последовательностей $\ell^\infty$:} $d(x, y) = \sup_{n \geqslant 1} |x_n - y_n|$.
	
	\begin{proof}
		Неравенство треугольника: $|x_n - z_n| \leqslant |x_n - y_n| + |y_n - z_n| \leqslant d(x,y) + d(y,z)$ для каждого $n$. Берём $\sup$.
	\end{proof}
	
	\textbf{6. Пространство непрерывных функций $C[a,b]$:} $d(f, g) = \max_{x \in [a,b]} |f(x) - g(x)|$.
	
	\begin{proof}
		1. $d(f,g) = 0 \Leftrightarrow |f(x) - g(x)| = 0$ для всех $x \Leftrightarrow f = g$.
		
		2. Симметричность: $|f(x) - g(x)| = |g(x) - f(x)|$.
		
		3. $|f(x) - h(x)| \leqslant |f(x) - g(x)| + |g(x) - h(x)| \leqslant d(f,g) + d(g,h)$. Берём максимум.
	\end{proof}
	
	\textbf{7. $p$-адическая метрика на $\mathbb{Z}$:} $d_p(a, b) = p^{-\nu_p(a-b)}$, где $\nu_p(m)$~--- показатель степени $p$ в разложении $m$ (при $m = 0$ полагают $d_p(a, a) = 0$).
	
	\begin{proof}
		Ключевое свойство: $\nu_p(a + b) \geqslant \min(\nu_p(a), \nu_p(b))$.
		
		Из него следует \textbf{ультраметрическое неравенство}: $d_p(a, c) \leqslant \max(d_p(a, b), d_p(b, c))$.
		
		$\nu_p(a - c) = \nu_p((a-b) + (b-c)) \geqslant \min(\nu_p(a-b), \nu_p(b-c))$, откуда $p^{-\nu_p(a-c)} \leqslant \max(p^{-\nu_p(a-b)}, p^{-\nu_p(b-c)})$.
		
		Ультраметрическое неравенство сильнее обычного: $\max(d(a,b), d(b,c)) \leqslant d(a,b) + d(b,c)$.
	\end{proof}
	
	\subsection{Открытые и замкнутые множества}
	
	\begin{definition}
		\textbf{Открытый шар:} $B(x, r) = \{y \in X : d(x, y) < r\}$.
		
		\textbf{Замкнутый шар:} $\overline{B}(x, r) = \{y \in X : d(x, y) \leqslant r\}$.
	\end{definition}
	
	\begin{definition}
		Множество $U \subseteq X$ \textbf{открыто}, если $\forall x \in U \ \exists r > 0: B(x, r) \subseteq U$.
		\end{definition}
	
	\begin{definition}
		Множество $F \subseteq X$ \textbf{замкнуто}, если $X \setminus F$ открыто.
	\end{definition}
	
	\begin{theorem}[Свойства открытых и замкнутых множеств]
		\begin{enumerate}
			\item $\varnothing$ и $X$ открыты (и замкнуты).
			\item Объединение любого семейства открытых множеств открыто.
			\item Пересечение конечного числа открытых множеств открыто.
			\item Пересечение любого семейства замкнутых множеств замкнуто.
			\item Объединение конечного числа замкнутых множеств замкнуто.
		\end{enumerate}
	\end{theorem}
	
	\begin{proof}
		1. $\varnothing$ открыто вакуумно. $X$ открыто: для любого $x \in X$ и любого $r > 0$: $B(x, r) \subseteq X$.
		
		2. Пусть $x \in \bigcup_\alpha U_\alpha$, тогда $x \in U_{\alpha_0}$ для некоторого $\alpha_0$. Так как $U_{\alpha_0}$ открыто, $\exists r > 0: B(x, r) \subseteq U_{\alpha_0} \subseteq \bigcup_\alpha U_\alpha$.
		
		3. Пусть $x \in U_1 \cap \dots \cap U_k$. Для каждого $i$: $\exists r_i > 0: B(x, r_i) \subseteq U_i$. Берём $r = \min(r_1, \dots, r_k) > 0$. Тогда $B(x, r) \subseteq U_i$ для всех $i$, значит $B(x, r) \subseteq U_1 \cap \dots \cap U_k$.
		
		4--5. Следует из 2--3 по де Моргану: дополнение пересечения замкнутых~--- объединение открытых, и наоборот.
	\end{proof}
	
	\begin{note}
		Бесконечное пересечение открытых может не быть открытым: $\bigcap_{n=1}^\infty (-1/n, 1/n) = \{0\}$~--- замкнуто в $\mathbb{R}$.
	\end{note}
	
	\begin{definition}
		\textbf{Замыкание} $\Cl A = \overline{A}$~--- наименьшее замкнутое множество, содержащее $A$:
		\[\Cl A = \bigcap \{F \subseteq X : F \text{ замкнуто}, A \subseteq F\}.\]
		
		Эквивалентное описание: $x \in \Cl A \Leftrightarrow \forall r > 0: B(x, r) \cap A \neq \varnothing$.
	\end{definition}
	
	\begin{definition}
		\textbf{Внутренность} $\operatorname{Int} A$~--- наибольшее открытое множество, содержащееся в $A$:
		\[\operatorname{Int} A = \bigcup \{U \subseteq X : U \text{ открыто}, U \subseteq A\}.\]
		
		Эквивалентное описание: $x \in \operatorname{Int} A \Leftrightarrow \exists r > 0: B(x, r) \subseteq A$.
	\end{definition}
	
	\subsection{Непрерывные отображения}
	
	\begin{definition}
		$f: X \to Y$ \textbf{непрерывно в точке} $x_0 \in X$, если:
		\[\forall \varepsilon > 0 \ \exists \delta > 0 \ \forall x \in X: d_X(x, x_0) < \delta \Rightarrow d_Y(f(x), f(x_0)) < \varepsilon.\]
		
		$f$ \textbf{непрерывно}, если оно непрерывно в каждой точке.
	\end{definition}
	
	\begin{theorem}[Равносильные определения непрерывности]
		Для отображения $f: X \to Y$ следующие условия равносильны:
		\begin{enumerate}
			\item $f$ непрерывно (в смысле $\varepsilon$-$\delta$ в каждой точке).
			\item Прообраз любого открытого множества открыт: $\forall U \subseteq Y$ открытого $f^{-1}(U)$ открыто в $X$.
			\item Прообраз любого замкнутого множества замкнут: $\forall F \subseteq Y$ замкнутого $f^{-1}(F)$ замкнуто в $X$.
			\item $f(\Cl A) \subseteq \Cl f(A)$ для любого $A \subseteq X$.
			\item Для любой сходящейся последовательности $x_n \to x$ в $X$: $f(x_n) \to f(x)$ в $Y$.
		\end{enumerate}
	\end{theorem}
	
	\begin{proof}
		$(1) \Rightarrow (2)$: Пусть $U$ открыто в $Y$, $x_0 \in f^{-1}(U)$. Тогда $f(x_0) \in U$, и $\exists \varepsilon > 0: B_Y(f(x_0), \varepsilon) \subseteq U$. По непрерывности в $x_0$: $\exists \delta > 0: f(B_X(x_0, \delta)) \subseteq B_Y(f(x_0), \varepsilon) \subseteq U$. Значит $B_X(x_0, \delta) \subseteq f^{-1}(U)$~--- $f^{-1}(U)$ открыто.
		
		$(2) \Rightarrow (1)$: Для $x_0 \in X$, $\varepsilon > 0$ рассмотрим $U = B_Y(f(x_0), \varepsilon)$~--- открытое. $f^{-1}(U)$ открыто и содержит $x_0$. Значит $\exists \delta > 0: B_X(x_0, \delta) \subseteq f^{-1}(U)$, то есть $d_X(x, x_0) < \delta \Rightarrow d_Y(f(x), f(x_0)) < \varepsilon$.
		
		$(2) \Leftrightarrow (3)$: $f^{-1}(Y \setminus F) = X \setminus f^{-1}(F)$. Дополнение открыто $\Leftrightarrow$ само замкнуто.
		
		$(1) \Rightarrow (5)$: Пусть $x_n \to x$, $\varepsilon > 0$. $\exists \delta > 0: d_X(y, x) < \delta \Rightarrow d_Y(f(y), f(x)) < \varepsilon$. $\exists N: n > N \Rightarrow d_X(x_n, x) < \delta$. Тогда $d_Y(f(x_n), f(x)) < \varepsilon$.
		
		$(5) \Rightarrow (4)$: Пусть $x \in \Cl A$. Тогда $\exists$ последовательность $x_n \in A: x_n \to x$. По (5): $f(x_n) \to f(x)$. Так как $f(x_n) \in f(A)$, имеем $f(x) \in \Cl f(A)$.
		
		$(4) \Rightarrow (3)$: Пусть $F$ замкнуто в $Y$, $A = f^{-1}(F)$. Тогда $f(A) \subseteq F$, значит $f(\Cl A) \subseteq \Cl f(A) \subseteq \Cl F = F$. Отсюда $\Cl A \subseteq f^{-1}(F) = A$, то есть $\Cl A = A$~--- $A$ замкнуто.
	\end{proof}
	
	\begin{note}
		Композиция непрерывных отображений непрерывна: если $f: X \to Y$ и $g: Y \to Z$ непрерывны, то $(g \circ f)^{-1}(U) = f^{-1}(g^{-1}(U))$~--- открыто как прообраз открытого при непрерывном $f$.
	\end{note}
	
	\subsection{Связность}
	
	\begin{definition}
		Метрическое пространство $X$ \textbf{связно}, если не существует разбиения $X = U \sqcup V$ на два непустых открытых множества.
		
		Подмножество $A \subseteq X$ \textbf{связно}, если $A$ связно как подпространство (с индуцированной метрикой $d|_{A \times A}$).
	\end{definition}
	
	\begin{note}
		Эквивалентно: $X$ связно $\Leftrightarrow$ единственные подмножества $X$, одновременно открытые и замкнутые~--- это $\varnothing$ и $X$.
	\end{note}
	
	\begin{theorem}
		Если $A \subseteq X$ связно, то $\Cl A$ тоже связно.
	\end{theorem}
	
	\begin{proof}
		Предположим $\Cl A = U \sqcup V$, где $U, V$ открыты в $\Cl A$ и непусты. Тогда $A = (A \cap U) \sqcup (A \cap V)$, где $A \cap U$ и $A \cap V$ открыты в $A$ (как пересечения $A$ с открытыми в $\Cl A$, а те открыты и в индуцированной метрике).
		
		Так как $A$ связно, либо $A \cap U = \varnothing$, либо $A \cap V = \varnothing$.
		
		Пусть $A \cap V = \varnothing$, тогда $A \subseteq U$. Множество $U$ открыто в $\Cl A$, значит его дополнение $V$ замкнуто в $\Cl A$. Так как $A \subseteq U \subseteq \Cl A \setminus V$ и $\Cl A \setminus V$ замкнуто в $\Cl A$, то $\Cl(A) \subseteq \Cl A \setminus V$, то есть $V = \varnothing$~--- противоречие.
	\end{proof}
	
	\begin{theorem}
		Непрерывный образ связного множества связен.
	\end{theorem}
	
	\begin{proof}
		Пусть $f: X \to Y$ непрерывно, $X$ связно. Предположим $f(X) = U \sqcup V$, где $U, V$ открыты в $f(X)$, непусты.
		
		Тогда $X = f^{-1}(U) \sqcup f^{-1}(V)$. По непрерывности $f$, прообразы открытых открыты в $X$. Оба непусты (т.к. $U, V \subseteq f(X)$). Противоречие со связностью $X$.
	\end{proof}
	
	\begin{theorem}
		Подмножество $\mathbb{R}$ связно $\Leftrightarrow$ оно является промежутком (то есть $\forall a, b \in A, a < b: [a, b] \subseteq A$).
	\end{theorem}
	
	\begin{proof}
		$(\Leftarrow)$ Пусть $I$~--- промежуток, $I = U \sqcup V$, $U, V$ открыты в $I$, непусты. Возьмём $a \in U$, $b \in V$, б.о.о. $a < b$.
		
		Пусть $c = \sup(U \cap [a, b])$. Тогда $c \in [a, b] \subseteq I$, и $c \in \Cl U$. Поскольку $V$ открыто и $V \cap U = \varnothing$, имеем $\Cl U \cap V = \varnothing$ (иначе точка из $\Cl U \cap V$ лежит в $V$ и является пределом точек из $U$~--- противоречие с открытостью $V$ и тем, что $V \cap U = \varnothing$). Значит $c \notin V$, то есть $c \in U$.
		
		Но $U$ открыто: $\exists \varepsilon > 0: (c - \varepsilon, c + \varepsilon) \cap I \subseteq U$. Тогда $c + \min(\varepsilon/2, b - c) \in U \cap [a, b]$ при $c < b$, противоречие с тем, что $c = \sup$. Значит $c = b$, но $b \in V$~--- противоречие с $c \in U$.
		
		$(\Rightarrow)$ Пусть $A$ связно, $a, b \in A$, $a < c < b$. Если $c \notin A$:
		$A = (A \cap (-\infty, c)) \sqcup (A \cap (c, +\infty))$~--- разбиение на непустые открытые в $A$. Противоречие.
	\end{proof}
	
	\begin{theorem}[Теорема о промежуточном значении]
		Пусть $f: [a, b] \to \mathbb{R}$ непрерывна, $f(a) < c < f(b)$ (или $f(a) > c > f(b)$). Тогда $\exists x \in (a, b): f(x) = c$.
	\end{theorem}
	
	\begin{proof}
		$[a, b]$ связен (как промежуток). $f([a, b])$ связен (непрерывный образ связного). Связное подмножество $\mathbb{R}$~--- промежуток. Так как $f(a)$ и $f(b)$ лежат в $f([a, b])$ и $c$ находится между ними, $c \in f([a, b])$, то есть $\exists x: f(x) = c$.
	\end{proof}
	
	\subsection{Задачи о блинах}
	
	\begin{theorem}[О разрезании двух фигур (Ham Sandwich в $\mathbb{R}^2$)]
		Любые две ограниченные измеримые фигуры на плоскости можно одновременно разрезать пополам одной прямой.
	\end{theorem}
	
	\begin{proof}[Идея доказательства через теорему о промежуточном значении]
		Параметризуем направления единичных векторов углом $\theta \in [0, \pi)$. Для каждого $\theta$ существует единственная прямая $\ell(\theta)$ с нормалью направления $\theta$, делящая первую фигуру $A_1$ ровно пополам (по площади).
		
		Пусть $f(\theta)$~--- разность площадей частей фигуры $A_2$, разделённых прямой $\ell(\theta)$: $f(\theta) = \text{площадь}(A_2 \text{ слева от } \ell(\theta)) - \text{площадь}(A_2 \text{ справа от } \ell(\theta))$.
		
		Функция $f$ непрерывна по $\theta$. При повороте на $\pi$ прямая $\ell$ та же, но «лево» и «право» меняются местами: $f(\theta + \pi) = -f(\theta)$ (считая $\theta$ по модулю $\pi$ с учётом ориентации).
		
		В частности, $f(0)$ и $f(\pi/2)$... нет, точнее: рассмотрим $f$ как функцию на $[0, \pi]$ с $f(\pi) = -f(0)$.
		
		Если $f(0) = 0$, готово. Иначе $f(0)$ и $f(\pi) = -f(0)$ имеют разные знаки. По теореме о промежуточном значении, $\exists \theta^* \in (0, \pi): f(\theta^*) = 0$~--- прямая $\ell(\theta^*)$ делит обе фигуры пополам.
	\end{proof}
	
	\begin{theorem}[Борсука--Улама для $S^1$]
		Для любой непрерывной функции $f: S^1 \to \mathbb{R}$ существует точка $x \in S^1$ такая, что $f(x) = f(-x)$.
	\end{theorem}
	
	\begin{proof}
		Рассмотрим $g: S^1 \to \mathbb{R}$, $g(x) = f(x) - f(-x)$. Тогда $g(-x) = f(-x) - f(x) = -g(x)$.
		
		Параметризуем $S^1$ углом $\theta \in [0, 2\pi]$. Тогда $g(0) + g(\pi) = g(\theta_0) - g(\theta_0) = 0$... нет. Точнее:
		
		$g(\pi) = f(-x_0) - f(x_0) = -g(0)$ (где $x_0$~--- точка, соответствующая $\theta = 0$, а $-x_0$~--- точка $\theta = \pi$).
		
		Если $g(0) = 0$, то $f(x_0) = f(-x_0)$~--- готово. Если $g(0) \neq 0$, то $g(0)$ и $g(\pi) = -g(0)$ разных знаков. По теореме о промежуточном значении (применяемой к пути от $0$ до $\pi$ на $S^1$): $\exists \theta^* \in (0, \pi): g(\theta^*) = 0$, то есть $f(x^*) = f(-x^*)$.
	\end{proof}
	
	\begin{note}
		Следствие теоремы Борсука--Улама: в любой момент времени на поверхности Земли существуют две антиподальные точки с одинаковой температурой и одинаковым атмосферным давлением.
	\end{note}
	
	\subsection{Компактность}
	
	\begin{definition}
		Семейство множеств $\{U_\alpha\}$ является \textbf{покрытием} множества $K$, если $K \subseteq \bigcup_\alpha U_\alpha$. Покрытие \textbf{открытое}, если все $U_\alpha$ открыты в $X$.
	\end{definition}
	
	\begin{definition}
		Множество $K \subseteq X$ \textbf{компактно}, если из любого открытого покрытия $K$ можно выбрать конечное подпокрытие.
	\end{definition}
	
	\begin{theorem}[Свойства компактных множеств]
		\begin{enumerate}
			\item Компактное подмножество метрического пространства замкнуто и ограничено.
			\item Замкнутое подмножество компактного множества компактно.
			\item Непрерывный образ компактного множества компактен.
		\end{enumerate}
	\end{theorem}
	
	\begin{proof}
		\textbf{1. Ограниченность:} зафиксируем $x_0 \in X$. Открытые шары $B(x_0, n)$, $n = 1, 2, \dots$, покрывают $X \supseteq K$. Выбираем конечное подпокрытие: $K \subseteq B(x_0, N)$ для некоторого $N$.
		
		\textbf{1. Замкнутость:} Пусть $y \notin K$. Для каждого $x \in K$ рассмотрим $r_x = d(x, y)/2 > 0$. Шары $B(x, r_x)$ покрывают $K$. Выбираем конечное подпокрытие: $K \subseteq B(x_1, r_1) \cup \dots \cup B(x_m, r_m)$.
		
		Пусть $r = \min_i r_i > 0$. Тогда $B(y, r) \cap B(x_i, r_i) = \varnothing$ для всех $i$ (т.к. $r \leqslant r_i = d(x_i, y)/2$, и если $z \in B(y, r) \cap B(x_i, r_i)$, то $d(x_i, y) \leqslant d(x_i, z) + d(z, y) < r_i + r \leqslant 2r_i = d(x_i, y)$~--- противоречие).
		
		Значит $B(y, r) \cap K = \varnothing$, то есть $y$ не является точкой прикосновения $K$, откуда $K$ замкнуто.
		
		\textbf{2.} Пусть $F \subseteq K$ замкнуто, $\{U_\alpha\}$~--- открытое покрытие $F$. Тогда $\{U_\alpha\} \cup \{X \setminus F\}$~--- открытое покрытие $K$. Выбираем конечное подпокрытие $K$; убираем $X \setminus F$~--- получаем конечное покрытие $F$.
		
		\textbf{3.} Пусть $f: K \to Y$ непрерывно, $\{V_\alpha\}$~--- открытое покрытие $f(K)$. Тогда $\{f^{-1}(V_\alpha)\}$~--- открытое покрытие $K$ (открытые по непрерывности). Выбираем конечное подпокрытие $f^{-1}(V_1), \dots, f^{-1}(V_m)$. Тогда $V_1, \dots, V_m$ покрывают $f(K)$.
	\end{proof}
	
	\begin{theorem}[Критерий компактности в $\mathbb{R}^n$]
		$K \subseteq \mathbb{R}^n$ компактно $\Leftrightarrow$ $K$ замкнуто и ограничено.
	\end{theorem}
	
	\begin{proof}
		$(\Rightarrow)$ Доказано в предыдущей теореме.
		
		$(\Leftarrow)$ Докажем, что замкнутый параллелепипед $P = [a_1, b_1] \times \dots \times [a_n, b_n]$ компактен (тогда $K \subseteq P$ для некоторого $P$, и как замкнутое подмножество компакта $K$ компактно).
		
		\textbf{Шаг 1: $[a, b]$ компактен в $\mathbb{R}$.}
		
		Пусть $\{U_\alpha\}$~--- открытое покрытие $[a, b]$ без конечного подпокрытия. Делим $[a, b]$ пополам: один из отрезков $[a, (a+b)/2]$ или $[(a+b)/2, b]$ тоже не имеет конечного подпокрытия (иначе объединением двух конечных покрывали бы $[a,b]$). Выбираем его и повторяем.
		
		Получаем вложенную последовательность отрезков $[a_k, b_k]$ длины $(b-a)/2^k$, каждый без конечного подпокрытия. По теореме о вложенных отрезках: $\exists c \in \bigcap_k [a_k, b_k]$. Тогда $c \in U_{\alpha_0}$ для некоторого $\alpha_0$, и $\exists \varepsilon > 0: (c - \varepsilon, c + \varepsilon) \subseteq U_{\alpha_0}$. При достаточно большом $k$: $[a_k, b_k] \subseteq (c-\varepsilon, c+\varepsilon) \subseteq U_{\alpha_0}$~--- конечное (из одного элемента) подпокрытие. Противоречие.
		
		\textbf{Шаг 2: прямое произведение компактов компактно.}
		
		Покажем для $P = K_1 \times K_2$ (индукция). Пусть $\{U_\alpha\}$ покрывает $P$. Для каждого $x \in K_1$ рассмотрим срез $\{x\} \times K_2 \cong K_2$~--- компакт. Выбираем конечное подпокрытие $U_{\alpha_1(x)}, \dots, U_{\alpha_k(x)}$.
		
		Для каждого $x$ существует окрестность $V_x \ni x$ (открытая в $K_1$) такая, что $V_x \times K_2 \subseteq U_{\alpha_1(x)} \cup \dots \cup U_{\alpha_k(x)}$ («лемма о трубке»: это следует из компактности $K_2$). Из $\{V_x\}_{x \in K_1}$~--- покрытия $K_1$~--- выбираем конечное: $V_{x_1}, \dots, V_{x_m}$. Тогда конечное семейство $\{U_{\alpha_j(x_i)}\}$ покрывает $P$.
		
		\textbf{Шаг 3:} $K$ замкнуто и ограничено $\Rightarrow$ $K \subseteq P$ для некоторого параллелепипеда $P$~--- компакта. По свойству 2 выше, $K$ компактно.
	\end{proof}
	
	\begin{theorem}[Секвенциальная компактность]
		В метрическом пространстве $K$ компактно $\Leftrightarrow$ из любой последовательности в $K$ можно выбрать подпоследовательность, сходящуюся к точке в $K$.
	\end{theorem}
	
	\begin{proof}[Для $\mathbb{R}^n$]
		$(\Rightarrow)$ Ограниченная последовательность в $\mathbb{R}^n$ (а $K$ ограничено): по теореме Больцано--Вейерштрасса содержит сходящуюся подпоследовательность. Предел лежит в $K$ (т.к. $K$ замкнуто).
		
		$(\Leftarrow)$ Если из любой последовательности можно выбрать сходящуюся в $K$ подпоследовательность, то $K$ замкнуто и ограничено~--- значит компактно в $\mathbb{R}^n$.
	\end{proof}
	
	\subsection{Теорема Вейерштрасса}
	
	\begin{theorem}[Вейерштрасс]
		Непрерывная функция на компакте ограничена и достигает своего максимума и минимума.
	\end{theorem}
	
	\begin{proof}
		Пусть $f: K \to \mathbb{R}$ непрерывна, $K$ компактно.
		
		\textbf{Ограниченность и компактность образа:} $f(K)$ компактно (непрерывный образ компакта), следовательно, замкнуто и ограничено в $\mathbb{R}$.
		
		\textbf{Достижение максимума:} Пусть $M = \sup f(K)$. Так как $f(K)$ ограничено, $M < +\infty$. Из определения супремума: $\exists$ последовательность $y_n \in f(K)$: $y_n \to M$.
		
		Так как $f(K)$ замкнуто: $M \in f(K)$, то есть $\exists x^* \in K: f(x^*) = M$.
		
		(Альтернативно: выберем $x_n \in K$ с $f(x_n) \to M$; по секвенциальной компактности $x_{n_k} \to x^* \in K$; по непрерывности $f(x^*) = \lim f(x_{n_k}) = M$.)
		
		Аналогично для минимума.
	\end{proof}
	
	\begin{theorem}[Равномерная непрерывность на компакте]
		Непрерывная функция на компакте равномерно непрерывна, то есть:
		\[\forall \varepsilon > 0 \ \exists \delta > 0 \ \forall x, y \in K: d(x, y) < \delta \Rightarrow |f(x) - f(y)| < \varepsilon.\]
	\end{theorem}
	
	\begin{proof}
		Предположим противное: $\exists \varepsilon_0 > 0$ такое, что $\forall \delta > 0 \ \exists x, y \in K: d(x, y) < \delta$, но $|f(x) - f(y)| \geqslant \varepsilon_0$.
		
		Берём $\delta_n = 1/n$: находим $x_n, y_n \in K$ с $d(x_n, y_n) < 1/n$ и $|f(x_n) - f(y_n)| \geqslant \varepsilon_0$.
		
		По секвенциальной компактности: $\exists$ подпоследовательность $x_{n_k} \to x^* \in K$. Тогда $d(y_{n_k}, x^*) \leqslant d(y_{n_k}, x_{n_k}) + d(x_{n_k}, x^*) < 1/n_k + d(x_{n_k}, x^*) \to 0$, то есть $y_{n_k} \to x^*$.
		
		По непрерывности $f$: $f(x_{n_k}) \to f(x^*)$ и $f(y_{n_k}) \to f(x^*)$.
		
		Но $|f(x_{n_k}) - f(y_{n_k})| \geqslant \varepsilon_0 > 0$~--- противоречие.
	\end{proof}

\end{document}